\documentclass[10pt,a4paper]{article}
\usepackage[margin=2.2cm]{geometry}
\usepackage{fontspec}
\setmonofont{DejaVu Sans Mono}[Scale=0.84]
\newfontfamily\unifont{DejaVu Serif}[Scale=0.9]
\usepackage{amsmath,amssymb,mathtools}
\setcounter{MaxMatrixCols}{20}
\usepackage{longtable,booktabs,array}
\usepackage{fancyvrb}
\usepackage[dvipsnames]{xcolor}
\usepackage[colorlinks=true,linkcolor=NavyBlue,urlcolor=NavyBlue]{hyperref}
\usepackage{microtype}
\usepackage{needspace}
\usepackage{tocloft}\setlength{\cftbeforesecskip}{0.45em}
\setlength{\parindent}{0pt}\setlength{\parskip}{0.55em}
\setlength{\emergencystretch}{8em}
\renewcommand{\arraystretch}{1.15}
\allowdisplaybreaks
\begin{document}
\title{Gravity on the 7+1 spacetime: provenance\\[0.4em]\large Sources, decisions, derivations and the verification record of the curvature, the Einstein and Lovelock tensors, the generalized frame matrices, the canonical spin connection, the covariant Dirac equation of dirac16complex and the three cosmological epochs — sections 11–14 of the notebook 7+1.nb}
\author{Repository: github.com/once-ere/7plus1 — prepared for Nash\\ Notebook: 7+1.nb (Wolfram Language 15.0.1)\\ Document generated by tools/make\_docs.py on 2026-10-27}
\date{}
\maketitle
\begin{abstract}
This record states what was asked, what was used, every decision taken with the equation behind it, how each result was derived and how it was checked: the 100 checks of sections 11–14 of 7+1.nb (100 True), the 101 boolean checks of the independent Python program tools/ref\_gravity.py (101 True), the comparison of the printed numbers, the program versions and the SHA-256 hashes of the files.
\end{abstract}
\tableofcontents
\Needspace{12\baselineskip}\section*{1. The task}\addcontentsline{toc}{section}{1. The task}
This stage of the project asked for the following (summarized; the original wording is the author’s private working material and is not reproduced):

\begin{enumerate}
\item Take the IDEAS of the repository github.com/once-ere/Dirac\_claude, which is cast on a 4+4 spacetime with four time-like coordinates and a definite warped metric, and rebuild them on the 7+1 spacetime of this project: ONE time-like coordinate $x_0$ and SEVEN space-like coordinates $x_1,\dots,x_7$. The 4+4 metric of Dirac\_claude is not used.
\item Einstein and Lovelock tensors of 7+1 spacetime.
\item Generalized frame matrices, a non-trivial canonical spin connection that couples gravity to dirac16complex, and a total covariant derivative under which the frame matrices are constant.
\item Kohn–Sham theory of dirac16complex in gravitational fields (primordial near $x_0=0$, modern, far future), its self-coupling, and the totality of the Standard Model (sections 15–19 of 7+1.nb; this record covers sections 11–14).
\item Compute the Lovelock tensors as in the author’s reference notebook Einstein-Lovelock-7+1-Nash.nb and do NOT use its section "FAILED Model metric g71; DO NOT USE!".
\item Make clear in all documentation, examples and provenance records that the xAct declaration $\{7,1,0\}$ declares seven space-like dimensions and one time-like dimension, while $\mathrm{diag}(+1,-1,\dots,-1)$ is simply the squares of the octonion basis (the author’s instruction of 27 October 2026).
\item For every documentation file a companion provenance file and a symbol file, each in Markdown and LaTeX, the LaTeX compiled to PDF; no document refers the reader to another one, and every equation used is written in the document that uses it.
\item Push every completed stage to github.com/once-ere/7plus1, respecting the repository’s .gitignore; use only the Wolfram front end WolframNB.exe of Wolfram 15.0.1.
\end{enumerate}
The coordinates of 7+1.nb are kept: $X=(x_0,\dots,x_7)$, $x_0$ the time (the direction of the octonion unit), $x_1,x_2,x_3$ ordinary space, $x_4,\dots,x_7$ auxiliary space.

\Needspace{8\baselineskip}\subsection*{1.1 The signature: the author’s explanation of 27 October 2026}\addcontentsline{toc}{subsection}{1.1 The signature: the author’s explanation of 27 October 2026}
\textbf{Source.} On 27 October 2026 the author of the project explained that the xAct list $\{7,1,0\}$ declares seven space-like dimensions and one time-like dimension, while $\mathrm{diag}(+1,-1,\dots,-1)$ is simply the squares of the octonion basis, and asked that this be made clear in all documentation, examples and provenance records. \textbf{What it corrected.} The version of 7+1.nb pushed in the commits 52d8729 to dc2f178 of github.com/once-ere/7plus1 declared the xAct metric of section 12 with \texttt{DefMetric[\{1, 7, 0\}, …]}, reading the list as the signs of $\eta$; since then section 12 declares \texttt{DefMetric[\{7, 1, 0\}, g71[-a, -b], CD, …]}, exactly as the author’s notebook Einstein-Lovelock-7+1-Nash.nb does. \textbf{What was added.} Checks in 7+1.nb — section 4.7: $\eta$ is the squares of the octonion basis, $e_Ae_A=\eta_{AA}e_0$ and $\mathrm{Re}(e_Ae_B)=\eta_{AB}$; section 12.2: \texttt{SignatureOfMetric[g71]} $=\{7,1,0\}$ with \texttt{SignDetOfMetric[g71]} $=-1$, $\mathrm{Re}(e_Ae_B)=\eta_{AB}$, and $\mathrm{Re}(x\,x)=\eta_{AB}x^Ax^B$; section 17: the flat metric $\mathrm{diag}(+1,-1,-1,-1)$ of the world brane is the squares of the basis $e_0,e_1,e_2,e_3$ of the quaternion subalgebra — and the same statements in the Python programs \texttt{tools/ref71.py}, \texttt{tools/ref\_gravity.py} and \texttt{tools/ref\_sm.py}; the explanation itself in the text of 7+1.nb (sections 1.2, 1.4, 4.2, 4.7, 7, 11, 12.2, 13, 15.1, 17) and in every document, symbol file and provenance file of the repository. \textbf{What did not change.} No result. Compared output by output with the backup below, every output of 7+1.nb is the same, except the declared list itself, the new checks and the check counts (and, as in every run, the name of the private context and the measured times); in particular every component of every Lovelock tensor is unchanged. The components follow from the component matrix of the metric, and the sign of the determinant, \texttt{SignDetOfMetric}, is $-1$ for both lists. \textbf{Backup.} The evaluated notebook as it was before is kept in the workspace as \texttt{backups/7+1\_2026-10-27\_before\_signature.nb} (the folder backups/ is excluded from the repository by .gitignore), 1034401 bytes, SHA-256 \texttt{060942488caf4f45120ffd37be0123e96992a31cd7fdbde09ee4dc4ced4f5d3a}; \texttt{tools/verify\_7p1\_nb.py} compares every input, text and section cell of 7+1.nb with it and finds every cell present, each one unchanged or changed by this explanation, and it compares every output in the same way.

\textbf{The explanation, in full.} The 7+1 spacetime has one time-like dimension ($x_0$, the direction of the octonion unit $e_0=1$) and seven space-like dimensions ($x_1,x_2,x_3$ ordinary space, $x_4,\dots,x_7$ auxiliary space). Two different objects describe this, and neither may be replaced by the other. (i) The tensor package xAct, used in section 12 of 7+1.nb exactly as in the author’s notebook Einstein-Lovelock-7+1-Nash.nb, declares the metric with \texttt{DefMetric[\{7, 1, 0\}, g71[-a, -b], CD, …]}: the list $\{7,1,0\}$ declares seven space-like dimensions, one time-like dimension and no degenerate dimension — the numbers of entries $+1$ and $-1$ of $\eta_{71}=\mathrm{diag}(-1,+1,\dots,+1)$ (\texttt{η71}), the same spacetime written in the sign convention of 4+4.nb (space-like $=+1$, time-like $=-1$). (ii) The matrix $\eta=-\eta_{71}=\mathrm{diag}(+1,-1,\dots,-1)$ (\texttt{η17}) is simply the squares of the octonion basis: $e_0^2=+1$ for the unit $e_0=1$ (time) and $e_i^2=-1$ for the seven imaginary units $e_1,\dots,e_7$ (space); more generally it is the real part of the product of two basis units, and it gives the quadratic form of the octonions:
\begin{gather*}
\eta_{AB}=\mathrm{Re}(e_Ae_B),\qquad e_0^2=+1,\qquad e_i^2=-1\ \ (i=1,\dots,7),\\ \mathrm{Re}(x\,x)=x_0^2-x_1^2-\cdots-x_7^2=\eta_{AB}\,x^Ax^B\qquad (x=x^Ae_A\in\mathbb O).
\end{gather*}
$\eta$ gives the signs of the Clifford relation $\Gamma_A\Gamma_B+\Gamma_B\Gamma_A=2\,\eta_{AB}\,1_{16}$ of the real 16×16 Dirac matrices: time squares to $+1_{16}$ and each of the seven space directions to $-1_{16}$, exactly like the octonion units (in the notation $Cl(p,q)$, with $p$ generators squaring to $+1$ and $q$ to $-1$, the algebra $Cl(1,7)\cong\mathbb R(16)$). It also gives the signs of the metric components: $g_{00}>0$ along the time-like direction, negative squares along the space-like directions. Every component computed by 7+1.nb follows from the component matrix of the metric; the declared list fixes for xAct the sign of the determinant of the metric, \texttt{SignDetOfMetric[g71]} $=-1$, which is the same in both sign conventions, since $\det \eta=\det \eta_{71}=-1$. 7+1.nb checks \texttt{SignatureOfMetric[g71]} $=\{7,1,0\}$, \texttt{SignDetOfMetric[g71]} $=-1$, $\mathrm{Re}(e_Ae_B)=\eta_{AB}$ and $\mathrm{Re}(x\,x)=\eta_{AB}x^Ax^B$.

\Needspace{12\baselineskip}\section*{2. Inputs}\addcontentsline{toc}{section}{2. Inputs}
\Needspace{8\baselineskip}\subsection*{2.1 Objects of the earlier sections of 7+1.nb}\addcontentsline{toc}{subsection}{2.1 Objects of the earlier sections of 7+1.nb}
Sections 11–14 use, unchanged, the coordinates \texttt{X = \{x0, …, x7\}}, the metric \texttt{η17} $=\eta$, the identity and zero matrices \texttt{ID16}, \texttt{Zero16}, the check registry \texttt{check[name, test]}, and the real 16×16 Dirac matrices \texttt{Γ17[A]} of section 4 with
\begin{equation*}
\Gamma_0=\begin{pmatrix}0&1_8\\1_8&0\end{pmatrix},\qquad \Gamma_i=\begin{pmatrix}\tau^O_i&0\\0&-\tau^O_i\end{pmatrix}\ (i=1,\dots,7),\qquad \Gamma_A\Gamma_B+\Gamma_B\Gamma_A=2\eta_{AB}1_{16},\qquad \Gamma^A=\eta^{AA}\Gamma_A,
\end{equation*}
where the real 8×8 matrices $\tau^O_i$ represent the seven imaginary octonion units ($(\tau^O_i)^2=-1_8$, $\tau^O_i\tau^O_j+\tau^O_j\tau^O_i=-2\delta_{ij}1_8$), and the Hamiltonian matrices \texttt{αD[j]} $=\Gamma^0\Gamma^j$, \texttt{βD} $=\Gamma_0$ of section 7. Before sections 11 and following were appended, the evaluated notebook was copied to \texttt{backups/7+1\_2026-10-26\_before\_gravity.nb} (excluded from the repository by .gitignore).

\begin{Verbatim}[fontsize=\small,frame=single,framesep=2mm]
backups/7+1_2026-10-26_before_gravity.nb
604743 bytes
SHA-256 438d972073e2fef9d7da6d57a541432b0fb1ea08424307f59ae2f3f7752bbfaf
\end{Verbatim}
\Needspace{8\baselineskip}\subsection*{2.2 The reference notebook Einstein-Lovelock-7+1-Nash.nb}\addcontentsline{toc}{subsection}{2.2 The reference notebook Einstein-Lovelock-7+1-Nash.nb}
The author’s notebook Einstein-Lovelock-7+1-Nash.nb (P. L. Nash, 2022, GNU GPL; corrected version of 27 October 2026, with its PDF printout) shows how the Lovelock tensors are to be computed. It is the author’s file and is not part of the repository. Taken from it, unchanged in substance:

\begin{itemize}
\item the xAct packages xCoba and xTras, the manifold \texttt{M8} with the abstract indices $a$, $b$, $c$, $d$, $f$, $i$, $j$, $k$, $l$, $m$, $n$, $p$, $q$, $r$, $s$, $y$, the chart \texttt{chartM8} with coordinate scalars \texttt{x0[]}, …, \texttt{x7[]}, and \texttt{MyArrayComponents};
\item the generalized Kronecker delta \texttt{kδ[lower\_, upper\_] := Det[Outer[delta, lower, upper]]};
\item the three Lovelock tensors with their normalization,
\end{itemize}
\begin{equation*}
\mathrm{Lovelock}_k=-\frac{1}{2^{k+1}}\,R_{c_1d_1}{}^{f_1i_1}\cdots R_{c_kd_k}{}^{f_ki_k}\;\delta^{b\,c_1d_1\cdots c_kd_k}_{a\,f_1i_1\cdots f_ki_k}\quad(k=1,2,3),\qquad \mathrm{Lovelock}_1=R_a{}^b-\tfrac12\delta_a{}^bR,
\end{equation*}
\begin{itemize}
\item the evaluation of components by assigning a component metric and setting \texttt{CD = LC[g]}, \texttt{RiemannCD = Riemann[CD]}, \texttt{RicciCD = Ricci[CD]}, \texttt{RicciScalarCD = RicciScalar[CD]}, \texttt{EinsteinCD = Einstein[CD]};
\item the form of the vacuum field equations and the procedure of solving them for $w_2,w_3$ (and $\Lambda$):
\end{itemize}
\begin{equation*}
0=-\Lambda\,\delta^\mu_\nu+H^{-2}w_1\,\mathrm{Lovelock}_1+H^{-4}w_2\,\mathrm{Lovelock}_2+H^{-6}w_3\,\mathrm{Lovelock}_3,
\end{equation*}
with the justification it quotes from D. Lovelock and H. Rund, \emph{Tensors, Differential Forms, and Variational Principles} (Dover), Remark 2 after Eq. (4.38): in $n$ dimensions the most general symmetric divergence-free tensor built from the metric and its first two derivatives is
\begin{equation*}
A^{lh}=\sqrt{|g|}\sum_{k=1}^{m-1}\alpha_{(k)}\,g^{jl}\,\delta^{h h_1\cdots h_{2k}}_{j j_1\cdots j_{2k}}\,R^{j_1j_2}{}_{h_1h_2}\cdots R^{j_{2k-1}j_{2k}}{}_{h_{2k-1}h_{2k}}+\lambda\sqrt{|g|}\,g^{lh},\qquad m-1=\frac82-1=3\ \ (n=8).
\end{equation*}
NOT taken: the section "FAILED Model metric g71; DO NOT USE!" of the reference notebook (its model metric and the solution found for it); its author states that it never worked for the 7+1 signature. No metric of that section appears in this project. Taken unchanged as well: the signature declaration \texttt{DefMetric[\{7, 1, 0\}, g71[-a, -b], CD, …]} — in xAct seven space-like dimensions ($x_1,\dots,x_7$), one time-like dimension ($x_0$), none degenerate — and the metric symbol \texttt{g71}.

\textbf{A correction (27 October 2026; section 1.1).} An earlier version of section 12 of 7+1.nb declared \texttt{DefMetric[\{1, 7, 0\}, …]}, reading the list as the signs of $\eta=\mathrm{diag}(+1,-1,\dots,-1)$; it now declares $\{7,1,0\}$, as the reference notebook does. The Lovelock tensors are the same with both declarations: they follow from the components of the metric, and \texttt{SignDetOfMetric} is $-1$ for both (a separate kernel with xTensor 1.2.0 confirmed this for a metric declared \texttt{\{1, 7, 0\}} and one declared \texttt{\{7, 1, 0\}}).

\begin{Verbatim}[fontsize=\small,frame=single,framesep=2mm]
Einstein-Lovelock-7+1-Nash.nb   SHA-256 b2dcb7b9f435d644da085fdd9e738a0cc80d3268ef28606ff9f367ed
    ...6828fdfc
Einstein-Lovelock-7+1-Nash.pdf  SHA-256 6580f9fcaff9c7f1337df35a7fce6b4b04c022d0e2d68954e8b3954b
    ...30f45125
\end{Verbatim}
\Needspace{8\baselineskip}\subsection*{2.3 Software and other sources}\addcontentsline{toc}{subsection}{2.3 Software and other sources}
\begin{itemize}
\item Wolfram Language 15.0.1 (the front end WolframNB.exe in \texttt{C:\textbackslash{}Program Files\textbackslash{}Wolfram Research\textbackslash{}Wolfram\textbackslash{}15.0.1}).
\item xAct (J. M. Martín-García, D. Yllanes, T. Nutma and others; GPL; www.xact.es), versions as printed by section 12.2 of 7+1.nb: xPerm 1.2.3, xTensor 1.2.0, xCoba 0.8.6, xTras 1.4.2.
\item Python 3.13.17 with sympy 1.14.0, numpy 2.5.3, scipy 1.18.1, mpmath 1.3.0 (the independent programs).
\item Ideas (no formulas or metrics) from the repository github.com/once-ere/Dirac\_claude: a 16-component field coupled to gravity by the canonical spin connection, Einstein–Lovelock gravity, a primordial pair-creation epoch, a Kohn–Sham treatment, the Standard Model on a lower-dimensional world surface.
\item The author’s private notes on vielbeins in $n$ dimensions (Cartan’s structure equation, a static metric with a radial coordinate, the Dirac equation, the Noether current and the U(1) coupling): their 7+1 versions were derived and checked anew (section 4 below gives the derivations, section 5 the checks). The notes are private and not part of the repository.
\item Standard textbook results, each derived and checked here: the Landau–Lifshitz curvature conventions, the Lanczos tensor, Kasner solutions, the Gibbons–Hawking temperature $H/(2\pi)$, the Landau–Zener formula.
\end{itemize}
\Needspace{12\baselineskip}\section*{3. Decisions}\addcontentsline{toc}{section}{3. Decisions}
{\footnotesize
\begin{longtable}{>{\raggedright\arraybackslash}p{\dimexpr 0.0550\linewidth-2\tabcolsep\relax}>{\raggedright\arraybackslash}p{\dimexpr 0.4950\linewidth-2\tabcolsep\relax}>{\raggedright\arraybackslash}p{\dimexpr 0.4500\linewidth-2\tabcolsep\relax}}
\toprule
\textbf{\#} & \textbf{decision} & \textbf{reason / equation} \\
\midrule\endfirsthead
\toprule
\textbf{\#} & \textbf{decision} & \textbf{reason / equation} \\
\midrule\endhead
\endfoot
\bottomrule\endlastfoot
1 & Christoffel symbols are written $\{{}^{\lambda}_{\mu\nu}\}$ & $\Gamma$ is reserved for the Dirac matrices \\
2 & Riemann tensor of Landau and Lifshitz: $R^\rho{}_{\sigma\mu\nu}=\partial_\mu\{{}^\rho_{\sigma\nu}\}-\partial_\nu\{{}^\rho_{\sigma\mu}\}+\{{}^\rho_{\mu\lambda}\}\{{}^\lambda_{\sigma\nu}\}-\{{}^\rho_{\nu\lambda}\}\{{}^\lambda_{\sigma\mu}\}$, $R_{\sigma\nu}=R^\rho{}_{\sigma\rho\nu}$ & with the signs of $\eta$ ($g_{00}>0$): $G^\mu{}_\nu-\bar\Lambda\delta^\mu_\nu=\kappa T^\mu{}_\nu$ with $T^0{}_0>0$; agrees with xAct (the reference notebook’s $\mathrm{Lovelock}_1$ equals this $G$) \\
3 & Lovelock tensors normalized as in the reference notebook, $E^{(k)}=-2^{-(k+1)}\delta RR\cdots$ & $E^{(1)}=G$, $E^{(2)}=$ Lanczos tensor \\
4 & two forms of the field equations: $w$-form (reference notebook) and $\alpha$-form & $\alpha_2=w_2/(w_1H^2)$, $\alpha_3=w_3/(w_1H^4)$, $\bar\Lambda=\Lambda H^2/w_1$ \\
5 & xAct signature declaration $\{7,1,0\}$, as in the reference notebook & seven space-like dimensions and one time-like dimension; $\eta=\mathrm{diag}(+1,-1,\dots,-1)$ is a different object, simply the squares of the octonion basis, $\eta_{AB}=\mathrm{Re}(e_Ae_B)$ \\
6 & xAct runs in the notebook kernel; the cell is guarded (\texttt{If[! MemberQ[\$Manifolds, M8], …]}, \texttt{Unset} of the component assignments); \texttt{tools/run\_eval.wl} evaluates the notebook in a fresh kernel & the cell can be evaluated twice while the notebook stays open; xAct keeps its definitions in the kernel, while a reopened notebook gets a new private context \\
7 & \texttt{xActLovelock} returns the Transpose of \texttt{MyArrayComponents} & the array has the lower index $a$ first; the pp-wave ($E^0{}_1=-E^1{}_0$) fixes the order \\
8 & a second, independent computation (pair-set sums) is kept & exact agreement with xAct is a strong check; it also works on exact rational tables \\
9 & the inhomogeneous metric is an explicit polynomial, checked exactly at $x_0=1/2$, $x_4=1/3$ & symbolic checks with arbitrary functions of two variables took hours \\
10 & exponential expansion is evaluated from the xAct components of the cosmological metric (\texttt{xlCos}), not by a new xAct run & xAct’s final Simplify with symbolic exponentials is very slow \\
11 & the spinor connection is built as \texttt{Zero16 + Sum[…]} & a sum of zeros would otherwise be the scalar 0, not a 16×16 matrix \\
12 & the 16-component mode equation is reduced with $\mathcal Z=\mathcal A\mathcal B\beta$, $\mathcal Z^2=-1$ & two classes of states ($\pm m\sigma_3$), created at different rates \\
13 & pair creation is integrated from $10^{-15}\lambda$ to $300\lambda$ (Kasner), from $-\ln2000$ to $\ln(10^6/M^2)$ (de Sitter) & adiabatic before and after; checked by $|\alpha|^2+|\beta|^2=1$ and the Landau–Zener formula \\*
\end{longtable}}
\Needspace{12\baselineskip}\section*{4. Derivations}\addcontentsline{toc}{section}{4. Derivations}
\textbf{Sectional curvatures of the cosmological metric.} In the frame $e=\mathrm{diag}(1,a,a,a,b,b,b,b)$ every $R^{AB}{}_{CD}$ is diagonal on index pairs, with
\begin{equation*}
K_{0i}=-(H_a'+H_a^2),\quad K_{0k}=-(H_b'+H_b^2),\quad K_{ij}=-H_a^2,\quad K_{kl}=-H_b^2,\quad K_{ik}=-H_aH_b,
\end{equation*}
In $L_k=2^{-k}\delta^{\alpha_1\beta_1\cdots}_{\gamma_1\delta_1\cdots}R^{\gamma_1\delta_1}{}_{\alpha_1\beta_1}\cdots$ every ordered sequence of $2k$ different indices contributes $2^k\prod_iK_{\alpha_i\beta_i}$ (each lower pair is the upper pair in one of its two orders), so $L_k$ is a sum over SETS of $k$ disjoint index pairs, and $E^{(k)\mu}{}_\mu$ (no sum) is $-\tfrac12$ times the same sum over the pairs that avoid $\mu$:
\begin{equation*}
L_k=2^kk!\sum_{\{p_1,\dots,p_k\}}K_{p_1}\cdots K_{p_k},\qquad E^{(k)\mu}{}_\mu=-2^{k-1}k!\sum_{\{p_1,\dots,p_k\}\not\ni\mu}K_{p_1}\cdots K_{p_k}.
\end{equation*}
Counting the pairs among $x_1,\dots,x_7$ (3 inside ordinary space, 6 inside the auxiliary space, 12 mixed) and the sets of two disjoint pairs (18 of type ordinary–auxiliary, 12 ordinary–mixed, 3 auxiliary–auxiliary, 36 auxiliary–mixed, 36 mixed–mixed) gives
\begin{gather*}
G^0{}_0=-\big(3K_{12}+6K_{45}+12K_{14}\big)=3H_a^2+12H_aH_b+6H_b^2,\\ E^{(2)0}{}_0=-4\big(12H_a^3H_b+54H_a^2H_b^2+36H_aH_b^3+3H_b^4\big)=-12H_b\big(4H_a^3+18H_a^2H_b+12H_aH_b^2+H_b^3\big).
\end{gather*}
\textbf{Kasner exponents.} With $a=x_0^p$, $b=x_0^q$: $H_a=p/x_0$, $H_a'=-p/x_0^2$ (likewise for $b$), and $x_0^{2k}E^{(k)}$ is a polynomial in $p,q$. Einstein: $3p+4q=1$, $3p^2+4q^2=1$. Gauss–Bonnet: $3p+4q=3$ and $-203p^3+81p^2+351p+27=0$. Cubic: $3p+4q=5$ and $(p,q)=(-15/7,20/7)$, besides the degenerate families $q=0$ and $p=0$.

\textbf{Spin connection of the cosmological metric.} From $\omega_\mu{}^A{}_B=e^A{}_\nu(\partial_\mu e_B{}^\nu+\{{}^\nu_{\mu\lambda}\}e_B{}^\lambda)$ with $\{{}^0_{ii}\}=aa'$, $\{{}^i_{0i}\}=a'/a$: $\omega_i{}^0{}_i=a'$, hence $\Omega_i=\tfrac14(\omega_{i0i}\Gamma^0\Gamma^i+\omega_{ii0}\Gamma^i\Gamma^0)=\tfrac{a'}{2}\Gamma_i\Gamma_0$.

\textbf{Static radial metric (Cartan).} $e^0=\sqrt f dx_0$, $e^1=\sqrt h dr$, $e^k=r dx^k$; $de^0=\frac{f'}{2\sqrt f}dr\wedge dx_0$, $de^k=dr\wedge dx^k$, $de^1=0$; Cartan’s equation $de^A+\omega^A{}_B\wedge e^B=0$ gives
\begin{equation*}
\omega^0{}_1=\frac{f'}{2\sqrt{fh}}\,dx_0,\quad \omega^k{}_1=\frac{dx^k}{\sqrt h},\quad \Omega_0=\frac{f'}{4\sqrt f\sqrt h}\Gamma_1\Gamma_0,\quad \Omega_k=\frac{1}{2\sqrt h}\Gamma_1\Gamma_k,\quad \Gamma^\mu\Omega_\mu=\frac{1}{\sqrt h}\Big(\frac{f'}{4f}+\frac3r\Big)\Gamma^1 .
\end{equation*}
\textbf{Exponential vacuum solutions.} For $a=e^{Hx_0}$, $b=e^{q_HHx_0}$ the vacuum equations are linear in $(w_2,w_3,\Lambda)$, and the divergence identity $\frac{d}{dx_0}E^{(k)0}{}_0+3H_a(E^{(k)0}{}_0-E^{(k)1}{}_1)+4H_b(E^{(k)0}{}_0-E^{(k)4}{}_4)=0$ with constant entries makes them of rank 2. For $q_H=-1$: $w_2=\Lambda/4-w_1/2$, $w_3=w_1/24-\Lambda/36$.

\textbf{de Sitter pair creation.} For $i\dot c=(ke^{-Hx_0}\sigma_1+M\sigma_3)c$ the change of variable $z\propto e^{-Hx_0}$ gives Whittaker functions; their connection formula yields $|\beta|^2=1/(e^{2\pi M/H}+1)$, checked numerically below.

\Needspace{12\baselineskip}\section*{5. Verification}\addcontentsline{toc}{section}{5. Verification}
\textbf{Evaluation of 7+1.nb.} The notebook was generated by \texttt{tools/make\_7p1\_nb.py} (which executes \texttt{tools/nb\_part\_gravity.py}, \texttt{tools/nb\_part\_ks\_gravity.py}, \texttt{tools/nb\_part\_sm.py}) and evaluated by \texttt{tools/run\_eval.wl} in a fresh kernel of Wolfram 15.0.1 for Microsoft Windows (64-bit) (July 2, 2026); seconds 271.569, message cells 0, finished Tue 27 Oct 2026 19:31:35. The final section of the notebook lists 399 checks; all passed: True.

\textbf{The checks of sections 11–14 of 7+1.nb:}

{\footnotesize
\begin{longtable}{>{\raggedright\arraybackslash}p{\dimexpr 0.0550\linewidth-2\tabcolsep\relax}>{\raggedright\arraybackslash}p{\dimexpr 0.7450\linewidth-2\tabcolsep\relax}>{\raggedright\arraybackslash}p{\dimexpr 0.0800\linewidth-2\tabcolsep\relax}>{\raggedright\arraybackslash}p{\dimexpr 0.1200\linewidth-2\tabcolsep\relax}}
\toprule
\textbf{\#} & \textbf{check} & \textbf{result} & \textbf{section} \\
\midrule\endfirsthead
\toprule
\textbf{\#} & \textbf{check} & \textbf{result} & \textbf{section} \\
\midrule\endhead
\endfoot
\bottomrule\endlastfoot
1 & \texttt{cosmological metric: sectional curvatures K\_0i = -(Ha'+Ha\textasciicircum{}2), K\_0k = -(Hb'+Hb\textasciicircum{}2), K\_ij = -Ha\textasciicircum{}2, K\_kl = -Hb\textasciicircum{}2, K\_ik = -Ha Hb} & True & 11.2 \\
2 & \texttt{cosmological metric: R\textasciicircum{}cd\_ab is non-zero only for \{c,d\} = \{a,b\}} & True & 11.2 \\
3 & \texttt{Einstein tensor: G\textasciicircum{}0\_0 = 3Ha\textasciicircum{}2 + 12 Ha Hb + 6 Hb\textasciicircum{}2, G\textasciicircum{}1\_1 = 2Ha' + 3Ha\textasciicircum{}2 + 4Hb' + 10Hb\textasciicircum{}2 + 8HaHb, G\textasciicircum{}4\_4 = 3Ha' + 6Ha\textasciicircum{}2 + 3Hb' + 6Hb\textasciicircum{}2 + 9HaHb} & True & 11.2 \\
4 & \texttt{scalar curvature R = -2 (3Ha' + 6Ha\textasciicircum{}2 + 4Hb' + 10Hb\textasciicircum{}2 + 12 Ha Hb)} & True & 11.2 \\
5 & \texttt{inhomogeneous metric, exactly at (x0, x4) = (1/2, 1/3): R\_rsmn = -R\_srmn = -R\_rsnm = R\_mnrs and R\textasciicircum{}r\_[smn] = 0 (first Bianchi identity)} & True & 11.3 \\
6 & \texttt{inhomogeneous metric: contracted Bianchi identity nabla\_mu G\textasciicircum{}mu\_nu == 0, exactly at (x0, x4) = (1/2, 1/3)} & True & 11.3 \\
7 & \texttt{cosmological metric: contracted Bianchi identity nabla\_mu G\textasciicircum{}mu\_nu = 0 (arbitrary a, b)} & True & 11.3 \\
8 & \texttt{generalized Kronecker delta: contraction identity delta\textasciicircum{}\{m1..mp r\}\_\{n1..np r\} = (8-p) delta\textasciicircum{}\{m1..mp\}\_\{n1..np\} (p = 1..4, 60 random index sets each)} & True & 12.1 \\
9 & \texttt{generalized Kronecker delta with 9 upper indices vanishes in 8 dimensions (100 random index sets)} & True & 12.1 \\
10 & \texttt{xAct declaration as in the reference notebook: SignatureOfMetric[g71] == \{7, 1, 0\} (seven space-like dimensions, one time-like dimension, none degenerate) and SignDetOfMetric[g71] == -1} & True & 12.2 \\
11 & \texttt{eta17 is simply the squares of the octonion basis: eta17[[A+1, B+1]] == Re(e\_A e\_B), e\_0\textasciicircum{}2 == +1 (time x0), e\_i\textasciicircum{}2 == -1 (space x\_i, i = 1..7)} & True & 12.2 \\
12 & \texttt{the quadratic form of the octonions has the signs of the squares of the octonion basis: Re(x x) == x0\textasciicircum{}2 - x1\textasciicircum{}2 - ... - x7\textasciicircum{}2 == x.eta17.x for x = x\textasciicircum{}A e\_A (symbolic)} & True & 12.2 \\
13 & \texttt{xAct: Lovelock1 == R\_a\textasciicircum{}b - (1/2) delta\_a\textasciicircum{}b R, the Einstein tensor (abstract indices)} & True & 12.3 \\
14 & \texttt{xAct: trace of Lovelock2 == -2 L\_2 with L\_2 = R\textasciicircum{}2 - 4 R\_cd R\textasciicircum{}cd + R\_cdfi R\textasciicircum{}cdfi (abstract indices)} & True & 12.3 \\
15 & \texttt{xAct, cosmological metric: Lovelock1 == EinsteinCD == the Einstein tensor of the curvature engine of section 11.1} & True & 12.4 \\
16 & \texttt{pair-set algorithm == brute-force generalized-delta contraction for E\textasciicircum{}(1) (cosmological metric)} & True & 12.5 \\
17 & \texttt{pair-set algorithm == xAct (section 12.4) for E\textasciicircum{}(1), E\textasciicircum{}(2), E\textasciicircum{}(3) of the cosmological metric (two independent computations)} & True & 12.5 \\
18 & \texttt{E\textasciicircum{}(1) == Einstein tensor G and L\_1 == R (cosmological metric)} & True & 12.6 \\
19 & \texttt{trace identity E\textasciicircum{}(k)mu\_mu == -((8-2k)/2) L\_k for k = 0..4} & True & 12.6 \\
20 & \texttt{E\textasciicircum{}(4) == 0 identically in 8 dimensions, while the Euler density L\_4 != 0} & True & 12.6 \\
21 & \texttt{every E\textasciicircum{}(k) is diagonal, E\textasciicircum{}(k)1\_1 = E\textasciicircum{}(k)2\_2 = E\textasciicircum{}(k)3\_3 and E\textasciicircum{}(k)4\_4 = ... = E\textasciicircum{}(k)7\_7} & True & 12.6 \\
22 & \texttt{E\textasciicircum{}(2) == Lanczos tensor and L\_2 == R\textasciicircum{}2 - 4 Ric.Ric + Riem.Riem (cosmological metric)} & True & 12.6 \\
23 & \texttt{E\textasciicircum{}(2)0\_0 = -12 Hb (4Ha\textasciicircum{}3 + 18Ha\textasciicircum{}2 Hb + 12 Ha Hb\textasciicircum{}2 + Hb\textasciicircum{}3) and E\textasciicircum{}(3)0\_0 = 360 Ha\textasciicircum{}2 Hb\textasciicircum{}3 (4Ha + 3Hb)} & True & 12.6 \\
24 & \texttt{divergence-free: d0 E\textasciicircum{}0\_0 + 3Ha (E\textasciicircum{}0\_0 - E\textasciicircum{}1\_1) + 4Hb (E\textasciicircum{}0\_0 - E\textasciicircum{}4\_4) == 0 for k = 1, 2, 3 (arbitrary a, b)} & True & 12.6 \\
25 & \texttt{variation of nL a\textasciicircum{}3 b\textasciicircum{}4 L\_k: EL\_nL = -2 a\textasciicircum{}3 b\textasciicircum{}4 E\textasciicircum{}(k)0\_0, EL\_a = -6 a\textasciicircum{}2 b\textasciicircum{}4 E\textasciicircum{}(k)1\_1, EL\_b = -8 a\textasciicircum{}3 b\textasciicircum{}3 E\textasciicircum{}(k)4\_4 (k = 1..4)} & True & 12.6 \\
26 & \texttt{inhomogeneous metric: nabla\_mu E\textasciicircum{}(k)mu\_nu == 0 EXACTLY at (x0, x4) = (1/2, 1/3) for k = 1, 2, 3} & True & 12.7 \\
27 & \texttt{inhomogeneous metric: E\textasciicircum{}(1) == G at the point, and E\textasciicircum{}(k)\_mu nu symmetric (k = 1, 2, 3)} & True & 12.7 \\
28 & \texttt{inhomogeneous metric: E\textasciicircum{}(2) == Lanczos tensor at the point (exact)} & True & 12.7 \\
29 & \texttt{pp-wave: E\textasciicircum{}(1)\_mu nu == +(1/2)(d2\textasciicircum{}2 + d3\textasciicircum{}2) fPP l\_mu l\_nu} & True & 12.8 \\
30 & \texttt{pp-wave: E\textasciicircum{}(2) == E\textasciicircum{}(3) == 0 for every profile fPP} & True & 12.8 \\
31 & \texttt{pp-wave: xAct (section 12.4) == pair-set algorithm for E\textasciicircum{}(1), E\textasciicircum{}(2), E\textasciicircum{}(3) (a metric that is not diagonal; arbitrary profile fPP)} & True & 12.8 \\
32 & \texttt{exponential expansion a = e\textasciicircum{}\{H x0\}, b = e\textasciicircum{}\{qH H x0\}: xAct Lovelock\_k == the closed forms of section 12.6 with Ha = H, Hb = qH H, Ha' = Hb' = 0 (k = 1, 2, 3)} & True & 12.9 \\
33 & \texttt{exponential expansion: three distinct vacuum equations, related by 3 Ha (eq00 - eq11) + 4 Hb (eq00 - eq44) == 0} & True & 12.9 \\
34 & \texttt{exponential vacuum solutions: (w2, w3) solve all 64 components for every qH (generic qH) and every Lambda} & True & 12.9 \\
35 & \texttt{eternal inflation (a = e\textasciicircum{}\{H x0\}) with deflation (b = e\textasciicircum{}\{-H x0\}): w2 == Lambda/4 - w1/2, w3 == w1/24 - Lambda/36; for Lambda = 0: w2/w1 == -1/2, w3/w1 == 1/24} & True & 12.9 \\
36 & \texttt{cosmological metric: Omega\_0 == 0, Omega\_i == (a'/2) Gamma17[i].Gamma17[0], Omega\_k == (b'/2) Gamma17[k].Gamma17[0]} & True & 13.2 \\
37 & \texttt{cosmological metric: Gamma\textasciicircum{}mu Omega\_mu == (3a'/(2a) + 2b'/b) Gamma17[0]} & True & 13.2 \\
38 & \texttt{current identity: d\_mu(sqrt|g| psibar Gamma\textasciicircum{}mu psi) == sqrt|g| [(D\_mu psibar) Gamma\textasciicircum{}mu psi + psibar Gamma\textasciicircum{}mu D\_mu psi] (psi, psibar arbitrary)} & True & 13.2 \\
39 & \texttt{Lichnerowicz formula: (Gamma\textasciicircum{}mu D\_mu)\textasciicircum{}2 psi == (g\textasciicircum{}mu nu nabla\_mu D\_nu - R/4) psi} & True & 13.2 \\
40 & \texttt{rescaling psi = (a\textasciicircum{}3 b\textasciicircum{}4)\textasciicircum{}(-1/2) chi removes the spin connection: Gamma\textasciicircum{}mu D\_mu psi == (a\textasciicircum{}3 b\textasciicircum{}4)\textasciicircum{}(-1/2) Gamma\textasciicircum{}mu d\_mu chi} & True & 13.2 \\
41 & \texttt{mode equation: Gamma17[0].(i Gamma\textasciicircum{}mu D\_mu - m) psi == (a\textasciicircum{}3 b\textasciicircum{}4)\textasciicircum{}(-1/2) e\textasciicircum{}\{i(p.x+q.y)\} (i dc/dx0 - h(x0) c) for psi = (a\textasciicircum{}3 b\textasciicircum{}4)\textasciicircum{}(-1/2) e\textasciicircum{}\{i(p.x+q.y)\} c(x0)} & True & 13.2 \\
42 & \texttt{conformal weight 7/2: Gamma\textasciicircum{}mu D\_mu (C\textasciicircum{}(-7/2) f) == C\textasciicircum{}(-9/2) Gamma17up[A] d\_A f for g = C(x0)\textasciicircum{}2 eta17} & True & 13.4 \\
43 & \texttt{Cartan's first structure equation de\textasciicircum{}A + omega\textasciicircum{}A\_B \textasciicircum{} e\textasciicircum{}B == 0 (no torsion): cosmological, pp-wave and inhomogeneous frames} & True & 13.5 \\
44 & \texttt{structure coefficients: omega\_C,AB == -(1/2)(C\_ACB + C\_BAC - C\_CBA) with C\textasciicircum{}A\_BC = e\_B\textasciicircum{}mu e\_C\textasciicircum{}nu (d\_mu e\textasciicircum{}A\_nu - d\_nu e\textasciicircum{}A\_mu): cosmological and inhomogeneous frames} & True & 13.5 \\
45 & \texttt{static radial metric: Cartan's first structure equation and omega\_C,AB == -(1/2)(C\_ACB + C\_BAC - C\_CBA)} & True & 13.5 \\
46 & \texttt{static radial metric: Omega\_0 == (f'/(4 Sqrt[f] Sqrt[h])) Gamma17[1].Gamma17[0], Omega\_1 == 0, Omega\_k == (1/(2 Sqrt[h])) Gamma17[1].Gamma17[k] (k = 2..7)} & True & 13.5 \\
47 & \texttt{static radial metric: Gamma\textasciicircum{}mu Omega\_mu == (1/Sqrt[h]) (f'/(4 f) + 3/r) Gamma17up[1]} & True & 13.5 \\
48 & \texttt{static radial metric: psi = f\textasciicircum{}(-1/4) r\textasciicircum{}(-3) chi removes the spin connection, Gamma\textasciicircum{}mu D\_mu psi == f\textasciicircum{}(-1/4) r\textasciicircum{}(-3) Gamma\textasciicircum{}mu d\_mu chi} & True & 13.5 \\
49 & \texttt{U(1) gauge covariance: Gamma\textasciicircum{}mu (D\_mu - i q A'\_mu)(e\textasciicircum{}\{i q alpha\} psi) == e\textasciicircum{}\{i q alpha\} Gamma\textasciicircum{}mu (D\_mu - i q A\_mu) psi with A'\_mu = A\_mu + d\_mu alpha} & True & 13.5 \\
50 & \texttt{current identity with the gauge field: d\_mu(sqrt|g| psibar Gamma\textasciicircum{}mu psi) == sqrt|g| [(D\_mu psibar) Gamma\textasciicircum{}mu psi + psibar Gamma\textasciicircum{}mu D\_mu psi], D including -i q A\_mu} & True & 13.5 \\
51 & \texttt{Kasner metric: the Einstein tensor computed directly from the metric equals the closed form (x0\textasciicircum{}2 G\textasciicircum{}mu\_mu, mu = 0, 1, 4)} & True & 14.1 \\
52 & \texttt{Einstein Kasner: (p, q) = ((1+2Sqrt[2])/7, (2-3Sqrt[2])/14) and ((1-2Sqrt[2])/7, (2+3Sqrt[2])/14) solve E\textasciicircum{}(1) == 0} & True & 14.1 \\
53 & \texttt{Gauss-Bonnet Kasner: 3 real roots of the cubic; (p, (3-3p)/4) solve E\textasciicircum{}(2) == 0} & True & 14.1 \\
54 & \texttt{cubic Lovelock Kasner: (p, q) = (-15/7, 20/7) solves E\textasciicircum{}(3) == 0; E\textasciicircum{}(3) == 0 identically for q = 0 and for p = 0} & True & 14.1 \\
55 & \texttt{2x2 reduction: A, B, betaD square to ID16 and anticommute; Zc = A.B.betaD: Zc.Zc == -ID16, Transpose[Zc] == -Zc, Zc commutes with A, B, betaD} & True & 14.2 \\
56 & \texttt{Kasner pair creation: |alpha|\textasciicircum{}2 + |beta|\textasciicircum{}2 == 1 (to 10\textasciicircum{}-8) for (|p|, |q|) = (1,1), (2,1), (1,3)} & True & 14.3 \\
57 & \texttt{Kasner scaling: (|p|, |q|) = (1.4302.., 2.5) has K = 1 and the same |beta|\textasciicircum{}2 as (1, 1) (to 10\textasciicircum{}-6)} & True & 14.3 \\
58 & \texttt{16x16 evolution == 4 copies of the 2x2 evolution with +m sigma3 and 4 with -m sigma3 (Kasner, m = 0.3, |p| = 1, |q| = 1.5, x0 from 0.05 to 40): average |beta|\textasciicircum{}2 agrees to 10\textasciicircum{}-8; unitarity; the two classes differ} & True & 14.3 \\
59 & \texttt{isotropic expansion, massless: no pair creation (|beta|\textasciicircum{}2 < 10\textasciicircum{}-20)} & True & 14.3 \\
60 & \texttt{static auxiliary space (Hb = 0): E\textasciicircum{}(1) = (3Ha\textasciicircum{}2, 2Ha' + 3Ha\textasciicircum{}2, 3Ha' + 6Ha\textasciicircum{}2), E\textasciicircum{}(2) = (0, 0, -12 Ha\textasciicircum{}2 (Ha' + Ha\textasciicircum{}2)), E\textasciicircum{}(3) = 0} & True & 14.4 \\
61 & \texttt{8D de Sitter: E\textasciicircum{}(1) = 21 H\textasciicircum{}2, E\textasciicircum{}(2) = -420 H\textasciicircum{}4, E\textasciicircum{}(3) = 2520 H\textasciicircum{}6 (all diagonal entries)} & True & 14.5 \\
62 & \texttt{de Sitter over a static auxiliary space: E\textasciicircum{}(1) = (3, 3, 6) H\textasciicircum{}2, E\textasciicircum{}(2) = (0, 0, -12) H\textasciicircum{}4, E\textasciicircum{}(3) = 0} & True & 14.5 \\
63 & \texttt{de Sitter pair creation: |beta|\textasciicircum{}2 == 1/(Exp[2 Pi M/H] + 1) to 10\textasciicircum{}-7, M/H = 0.25, 0.5, 1} & True & 14.5 \\
64 & \texttt{integrator benchmark (Landau-Zener): |beta|\textasciicircum{}2 == Exp[-Pi kappa\textasciicircum{}2] to 10\textasciicircum{}-6} & True & 14.5 \\
65 & \texttt{cosmological metric: eta17 e e == g and \{Gamma\_mu, Gamma\_nu\} == 2 g\_mu nu ID16 (generalized frame matrices, real)} & True & 13.1 \\
66 & \texttt{cosmological metric: canonical spin connection antisymmetric, omega\_mu AB == -omega\_mu BA} & True & 13.1 \\
67 & \texttt{cosmological metric: vielbein postulate d\_mu e\textasciicircum{}A\_nu - \{l mu nu\} e\textasciicircum{}A\_l + omega\_mu\textasciicircum{}A\_B e\textasciicircum{}B\_nu == 0} & True & 13.1 \\
68 & \texttt{cosmological metric: TOTAL covariant derivative of the frame matrices vanishes, d\_mu Gamma\_nu - \{l mu nu\} Gamma\_l + [Omega\_mu, Gamma\_nu] == 0} & True & 13.1 \\
69 & \texttt{cosmological metric: constant frame matrices are covariantly constant, omega\_mu\textasciicircum{}A\_B Gamma17up[B] + [Omega\_mu, Gamma17up[A]] == 0} & True & 13.1 \\
70 & \texttt{cosmological metric: Omega\_mu real and Gamma17[0].Transpose[Omega\_mu].Gamma17[0] == -Omega\_mu (so D\_mu psibar = d\_mu psibar - psibar Omega\_mu)} & True & 13.1 \\
71 & \texttt{cosmological metric: spin curvature F\_mu nu == (1/4) R\_AB mu nu Gamma17up[A].Gamma17up[B]} & True & 13.1 \\
72 & \texttt{cosmological metric: Bianchi identity d\_l F\_mn + [Omega\_l, F\_mn] + cyclic == 0} & True & 13.1 \\
73 & \texttt{cosmological metric: (1/2) Gamma\textasciicircum{}mu Gamma\textasciicircum{}nu F\_mu nu == -(R/4) ID16} & True & 13.1 \\
74 & \texttt{pp-wave: eta17 e e == g and \{Gamma\_mu, Gamma\_nu\} == 2 g\_mu nu ID16 (generalized frame matrices, real)} & True & 13.1 \\
75 & \texttt{pp-wave: canonical spin connection antisymmetric, omega\_mu AB == -omega\_mu BA} & True & 13.1 \\
76 & \texttt{pp-wave: vielbein postulate d\_mu e\textasciicircum{}A\_nu - \{l mu nu\} e\textasciicircum{}A\_l + omega\_mu\textasciicircum{}A\_B e\textasciicircum{}B\_nu == 0} & True & 13.1 \\
77 & \texttt{pp-wave: TOTAL covariant derivative of the frame matrices vanishes, d\_mu Gamma\_nu - \{l mu nu\} Gamma\_l + [Omega\_mu, Gamma\_nu] == 0} & True & 13.1 \\
78 & \texttt{pp-wave: constant frame matrices are covariantly constant, omega\_mu\textasciicircum{}A\_B Gamma17up[B] + [Omega\_mu, Gamma17up[A]] == 0} & True & 13.1 \\
79 & \texttt{pp-wave: Omega\_mu real and Gamma17[0].Transpose[Omega\_mu].Gamma17[0] == -Omega\_mu (so D\_mu psibar = d\_mu psibar - psibar Omega\_mu)} & True & 13.1 \\
80 & \texttt{pp-wave: spin curvature F\_mu nu == (1/4) R\_AB mu nu Gamma17up[A].Gamma17up[B]} & True & 13.1 \\
81 & \texttt{pp-wave: Bianchi identity d\_l F\_mn + [Omega\_l, F\_mn] + cyclic == 0} & True & 13.1 \\
82 & \texttt{pp-wave: (1/2) Gamma\textasciicircum{}mu Gamma\textasciicircum{}nu F\_mu nu == -(R/4) ID16} & True & 13.1 \\
83 & \texttt{inhomogeneous metric (exact at (x0, x4) = (1/2, 1/3)): eta17 e e == g and \{Gamma\_mu, Gamma\_nu\} == 2 g\_mu nu ID16 (generalized frame matrices, real)} & True & 13.1 \\
84 & \texttt{inhomogeneous metric (exact at (x0, x4) = (1/2, 1/3)): canonical spin connection antisymmetric, omega\_mu AB == -omega\_mu BA} & True & 13.1 \\
85 & \texttt{inhomogeneous metric (exact at (x0, x4) = (1/2, 1/3)): vielbein postulate d\_mu e\textasciicircum{}A\_nu - \{l mu nu\} e\textasciicircum{}A\_l + omega\_mu\textasciicircum{}A\_B e\textasciicircum{}B\_nu == 0} & True & 13.1 \\
86 & \texttt{inhomogeneous metric (exact at (x0, x4) = (1/2, 1/3)): TOTAL covariant derivative of the frame matrices vanishes, d\_mu Gamma\_nu - \{l mu nu\} Gamma\_l + [Omega\_mu, Gamma\_nu] == 0} & True & 13.1 \\
87 & \texttt{inhomogeneous metric (exact at (x0, x4) = (1/2, 1/3)): constant frame matrices are covariantly constant, omega\_mu\textasciicircum{}A\_B Gamma17up[B] + [Omega\_mu, Gamma17up[A]] == 0} & True & 13.1 \\
88 & \texttt{inhomogeneous metric (exact at (x0, x4) = (1/2, 1/3)): Omega\_mu real and Gamma17[0].Transpose[Omega\_mu].Gamma17[0] == -Omega\_mu (so D\_mu psibar = d\_mu psibar - psibar Omega\_mu)} & True & 13.1 \\
89 & \texttt{inhomogeneous metric (exact at (x0, x4) = (1/2, 1/3)): spin curvature F\_mu nu == (1/4) R\_AB mu nu Gamma17up[A].Gamma17up[B]} & True & 13.1 \\
90 & \texttt{inhomogeneous metric (exact at (x0, x4) = (1/2, 1/3)): Bianchi identity d\_l F\_mn + [Omega\_l, F\_mn] + cyclic == 0} & True & 13.1 \\
91 & \texttt{inhomogeneous metric (exact at (x0, x4) = (1/2, 1/3)): (1/2) Gamma\textasciicircum{}mu Gamma\textasciicircum{}nu F\_mu nu == -(R/4) ID16} & True & 13.1 \\
92 & \texttt{static radial metric: eta17 e e == g and \{Gamma\_mu, Gamma\_nu\} == 2 g\_mu nu ID16 (generalized frame matrices, real)} & True & 13.1 \\
93 & \texttt{static radial metric: canonical spin connection antisymmetric, omega\_mu AB == -omega\_mu BA} & True & 13.1 \\
94 & \texttt{static radial metric: vielbein postulate d\_mu e\textasciicircum{}A\_nu - \{l mu nu\} e\textasciicircum{}A\_l + omega\_mu\textasciicircum{}A\_B e\textasciicircum{}B\_nu == 0} & True & 13.1 \\
95 & \texttt{static radial metric: TOTAL covariant derivative of the frame matrices vanishes, d\_mu Gamma\_nu - \{l mu nu\} Gamma\_l + [Omega\_mu, Gamma\_nu] == 0} & True & 13.1 \\
96 & \texttt{static radial metric: constant frame matrices are covariantly constant, omega\_mu\textasciicircum{}A\_B Gamma17up[B] + [Omega\_mu, Gamma17up[A]] == 0} & True & 13.1 \\
97 & \texttt{static radial metric: Omega\_mu real and Gamma17[0].Transpose[Omega\_mu].Gamma17[0] == -Omega\_mu (so D\_mu psibar = d\_mu psibar - psibar Omega\_mu)} & True & 13.1 \\
98 & \texttt{static radial metric: spin curvature F\_mu nu == (1/4) R\_AB mu nu Gamma17up[A].Gamma17up[B]} & True & 13.1 \\
99 & \texttt{static radial metric: Bianchi identity d\_l F\_mn + [Omega\_l, F\_mn] + cyclic == 0} & True & 13.1 \\
100 & \texttt{static radial metric: (1/2) Gamma\textasciicircum{}mu Gamma\textasciicircum{}nu F\_mu nu == -(R/4) ID16} & True & 13.1 \\*
\end{longtable}}
\textbf{The independent Python program} tools/ref\_gravity.py (its own curvature and Lovelock engine tools/lovelock71.py; the Dirac matrices rebuilt from the split octonions of 4+4.nb by tools/ref71.py). Its boolean checks:

{\footnotesize
\begin{longtable}{>{\raggedright\arraybackslash}p{\dimexpr 0.0550\linewidth-2\tabcolsep\relax}>{\raggedright\arraybackslash}p{\dimexpr 0.8250\linewidth-2\tabcolsep\relax}>{\raggedright\arraybackslash}p{\dimexpr 0.1200\linewidth-2\tabcolsep\relax}}
\toprule
\textbf{\#} & \textbf{check} & \textbf{result} \\
\midrule\endfirsthead
\toprule
\textbf{\#} & \textbf{check} & \textbf{result} \\
\midrule\endhead
\endfoot
\bottomrule\endlastfoot
1 & \texttt{eta17 is simply the squares of the octonion basis: eta\_AB = Re(e\_A e\_B), e\_0\textasciicircum{}2 = +1, e\_i\textasciicircum{}2 = -1 (i = 1..7)} & True \\
2 & \texttt{the quadratic form of the octonions has the signs of the squares of the octonion basis: Re(x x) = x0\textasciicircum{}2 - x1\textasciicircum{}2 - ... - x7\textasciicircum{}2 = x.eta17.x} & True \\
3 & \texttt{cosmological metric: sectional curvatures R\textasciicircum{}\{AB\}\_\{AB\} = -(Ha'+Ha\textasciicircum{}2), -(Hb'+Hb\textasciicircum{}2), -Ha\textasciicircum{}2, -Hb\textasciicircum{}2, -Ha Hb} & True \\
4 & \texttt{cosmological metric: Riemann R\textasciicircum{}\{cd\}\_\{ab\} non-zero only for \{c,d\} = \{a,b\}} & True \\
5 & \texttt{G\textasciicircum{}0\_0 = 3Ha\textasciicircum{}2 + 12 Ha Hb + 6 Hb\textasciicircum{}2} & True \\
6 & \texttt{G\textasciicircum{}1\_1 = 2Ha' + 3Ha\textasciicircum{}2 + 4Hb' + 10Hb\textasciicircum{}2 + 8HaHb} & True \\
7 & \texttt{G\textasciicircum{}4\_4 = 3Ha' + 6Ha\textasciicircum{}2 + 3Hb' + 6Hb\textasciicircum{}2 + 9HaHb} & True \\
8 & \texttt{scalar curvature R = -2(3Ha' + 6Ha\textasciicircum{}2 + 4Hb' + 10Hb\textasciicircum{}2 + 12 HaHb)} & True \\
9 & \texttt{E\textasciicircum{}(1) = Einstein tensor G (cosmological metric)} & True \\
10 & \texttt{L\_1 = R} & True \\
11 & \texttt{E\textasciicircum{}(0) = -delta/2} & True \\
12 & \texttt{trace identity E\textasciicircum{}(k)mu\_mu = -(8-2k)/2 L\_k, k = 0..4} & True \\
13 & \texttt{E\textasciicircum{}(4) = 0 identically in 8 dimensions, but L\_4 (Euler density) is not 0} & True \\
14 & \texttt{E\textasciicircum{}(k) diagonal: E\textasciicircum{}(k)mu\_nu = 0 for mu != nu, and E\textasciicircum{}(k)1\_1 = E\textasciicircum{}(k)2\_2 = E\textasciicircum{}(k)3\_3, E\textasciicircum{}(k)4\_4 = ... = E\textasciicircum{}(k)7\_7} & True \\
15 & \texttt{E\textasciicircum{}(2)0\_0 = -12 Hb (4Ha\textasciicircum{}3 + 18Ha\textasciicircum{}2Hb + 12HaHb\textasciicircum{}2 + Hb\textasciicircum{}3)} & True \\
16 & \texttt{E\textasciicircum{}(3)0\_0 = 360 Ha\textasciicircum{}2 Hb\textasciicircum{}3 (4Ha + 3Hb)} & True \\
17 & \texttt{E\textasciicircum{}(2) = Lanczos tensor 2(R R\_mn - 2R\_ma R\textasciicircum{}a\_n - 2R\textasciicircum{}ab R\_manb + R\_m\textasciicircum{}abc R\_nabc) - g\_mn L\_2/2, and L\_2 = R\textasciicircum{}2 - 4 Ric.Ric + Riem.Riem (cosmological metric)} & True \\
18 & \texttt{divergence-free E\textasciicircum{}(k), k = 1, 2, 3 (cosmological metric, arbitrary a, b)} & True \\
19 & \texttt{variation of sqrt|g| L\_k (lapse N, a, b): EL\_N = -2 a\textasciicircum{}3 b\textasciicircum{}4 E\textasciicircum{}(k)0\_0, EL\_a = -6 a\textasciicircum{}2 b\textasciicircum{}4 E\textasciicircum{}(k)1\_1, EL\_b = -8 a\textasciicircum{}3 b\textasciicircum{}3 E\textasciicircum{}(k)4\_4, k = 1..4} & True \\
20 & \texttt{inhomogeneous metric (N, A, B of x0 and x4): E\textasciicircum{}(1) = G at the point} & True \\
21 & \texttt{inhomogeneous metric: nabla\_mu E\textasciicircum{}(k)mu\_nu = 0 exactly at x0 = 1/2, x4 = 1/3, k = 1, 2, 3} & True \\
22 & \texttt{inhomogeneous metric: E\textasciicircum{}(k)\_mu nu symmetric, k = 1, 2, 3} & True \\
23 & \texttt{pp-wave: E\textasciicircum{}(1)\_mu nu = +(1/2) (d2\textasciicircum{}2 + d3\textasciicircum{}2) F l\_mu l\_nu, l = dx0 - dx1} & True \\
24 & \texttt{pp-wave: E\textasciicircum{}(2) = E\textasciicircum{}(3) = 0 for every profile F} & True \\
25 & \texttt{Einstein Kasner: p = (1+-2 sqrt2)/7, q = (2-+3 sqrt2)/14 solve E\textasciicircum{}(1) = 0; 3p+4q = 1, 3p\textasciicircum{}2+4q\textasciicircum{}2 = 1} & True \\
26 & \texttt{Gauss-Bonnet Kasner: the three (p, (3-3p)/4) solve E\textasciicircum{}(2) = 0, sum of the 7 exponents = 3} & True \\
27 & \texttt{cubic Lovelock Kasner: (p, q) = (-15/7, 20/7) solves E\textasciicircum{}(3) = 0, 3p + 4q = 5; degenerate families q = 0 and p = 0} & True \\
28 & \texttt{Einstein Kasner is not a cubic-Lovelock vacuum (E\textasciicircum{}(3) != 0)} & True \\
29 & \texttt{8D de Sitter (a = b = e\textasciicircum{}\{H x0\}): E\textasciicircum{}(1) = 21 H\textasciicircum{}2, E\textasciicircum{}(2) = -420 H\textasciicircum{}4, E\textasciicircum{}(3) = 2520 H\textasciicircum{}6 (times delta)} & True \\
30 & \texttt{4D de Sitter on a static T\textasciicircum{}4 (a = e\textasciicircum{}\{H x0\}, b = b0): E\textasciicircum{}(1) = diag(3,3,6) H\textasciicircum{}2, E\textasciicircum{}(2) = diag(0,0,-12) H\textasciicircum{}4, E\textasciicircum{}(3) = 0} & True \\
31 & \texttt{static T\textasciicircum{}4 (Hb = 0): E\textasciicircum{}(1) = (3Ha\textasciicircum{}2, 2Ha'+3Ha\textasciicircum{}2, 3Ha'+6Ha\textasciicircum{}2), E\textasciicircum{}(2) = (0, 0, -12 Ha\textasciicircum{}2 (Ha'+Ha\textasciicircum{}2)), E\textasciicircum{}(3) = 0} & True \\
32 & \texttt{exponential expansion a = e\textasciicircum{}\{H x0\}, b = e\textasciicircum{}\{q H x0\}: E\textasciicircum{}(k) (00, 11, 44) = closed forms with Ha = H, Hb = q H, Ha' = Hb' = 0 (k = 1, 2, 3)} & True \\
33 & \texttt{exponential expansion: 3 Ha (eq00 - eq11) + 4 Hb (eq00 - eq44) = 0 identically (the three vacuum equations have rank 2)} & True \\
34 & \texttt{exponential vacuum solutions: one solution (w2, w3) for given w1, Lambda, q, satisfying all three equations} & True \\
35 & \texttt{eternal inflation a = e\textasciicircum{}\{H x0\} with deflation b = e\textasciicircum{}\{-H x0\}: w2 = Lambda/4 - w1/2, w3 = w1/24 - Lambda/36 (Lambda = 0: w2/w1 = -1/2, w3/w1 = 1/24)} & True \\
36 & \texttt{cosmological metric: eta e e = g} & True \\
37 & \texttt{cosmological metric: \{Gamma\_mu, Gamma\_nu\} = 2 g\_mu nu} & True \\
38 & \texttt{cosmological metric: Gamma\_mu real} & True \\
39 & \texttt{cosmological metric: omega\_mu AB = -omega\_mu BA} & True \\
40 & \texttt{cosmological metric: vielbein postulate d\_mu e\textasciicircum{}A\_nu - chr\textasciicircum{}l\_mu nu e\textasciicircum{}A\_l + omega\_mu\textasciicircum{}A\_B e\textasciicircum{}B\_nu = 0} & True \\
41 & \texttt{cosmological metric: Omega\_mu real, Gamma17[0] Omega\_mu\textasciicircum{}T Gamma17[0] = -Omega\_mu} & True \\
42 & \texttt{cosmological metric: total covariant derivative of the frame matrices: d\_mu Gamma\_nu - chr\textasciicircum{}l\_mu nu Gamma\_l + [Omega\_mu, Gamma\_nu] = 0} & True \\
43 & \texttt{cosmological metric: constant frame matrices: omega\_mu\textasciicircum{}A\_B Gamma\textasciicircum{}B + [Omega\_mu, Gamma\textasciicircum{}A] = 0} & True \\
44 & \texttt{cosmological metric: spin curvature F\_mu nu = (1/4) R\_AB mu nu Gamma\textasciicircum{}A Gamma\textasciicircum{}B} & True \\
45 & \texttt{cosmological metric: Bianchi identity d\_l F\_mn + [Omega\_l, F\_mn] + cyclic = 0} & True \\
46 & \texttt{cosmological metric: (1/2) Gamma\textasciicircum{}mu Gamma\textasciicircum{}nu F\_mu nu = -(R/4) 1\_16} & True \\
47 & \texttt{cosmological metric: Gamma\textasciicircum{}mu Omega\_mu = (3a'/(2a) + 2b'/b) Gamma17[0]} & True \\
48 & \texttt{cosmological metric: Omega\_i = (a'/2) Gamma17[i] Gamma17[0] (i = 1..3), Omega\_k = (b'/2) Gamma17[k] Gamma17[0] (k = 4..7), Omega\_0 = 0} & True \\
49 & \texttt{pp-wave: eta e e = g} & True \\
50 & \texttt{pp-wave: \{Gamma\_mu, Gamma\_nu\} = 2 g\_mu nu} & True \\
51 & \texttt{pp-wave: Gamma\_mu real} & True \\
52 & \texttt{pp-wave: omega\_mu AB = -omega\_mu BA} & True \\
53 & \texttt{pp-wave: vielbein postulate d\_mu e\textasciicircum{}A\_nu - chr\textasciicircum{}l\_mu nu e\textasciicircum{}A\_l + omega\_mu\textasciicircum{}A\_B e\textasciicircum{}B\_nu = 0} & True \\
54 & \texttt{pp-wave: Omega\_mu real, Gamma17[0] Omega\_mu\textasciicircum{}T Gamma17[0] = -Omega\_mu} & True \\
55 & \texttt{pp-wave: total covariant derivative of the frame matrices: d\_mu Gamma\_nu - chr\textasciicircum{}l\_mu nu Gamma\_l + [Omega\_mu, Gamma\_nu] = 0} & True \\
56 & \texttt{pp-wave: constant frame matrices: omega\_mu\textasciicircum{}A\_B Gamma\textasciicircum{}B + [Omega\_mu, Gamma\textasciicircum{}A] = 0} & True \\
57 & \texttt{pp-wave: spin curvature F\_mu nu = (1/4) R\_AB mu nu Gamma\textasciicircum{}A Gamma\textasciicircum{}B} & True \\
58 & \texttt{pp-wave: Bianchi identity d\_l F\_mn + [Omega\_l, F\_mn] + cyclic = 0} & True \\
59 & \texttt{pp-wave: (1/2) Gamma\textasciicircum{}mu Gamma\textasciicircum{}nu F\_mu nu = -(R/4) 1\_16} & True \\
60 & \texttt{inhomogeneous metric N, A, B (x0, x4): eta e e = g} & True \\
61 & \texttt{inhomogeneous metric N, A, B (x0, x4): \{Gamma\_mu, Gamma\_nu\} = 2 g\_mu nu} & True \\
62 & \texttt{inhomogeneous metric N, A, B (x0, x4): Gamma\_mu real} & True \\
63 & \texttt{inhomogeneous metric N, A, B (x0, x4): omega\_mu AB = -omega\_mu BA} & True \\
64 & \texttt{inhomogeneous metric N, A, B (x0, x4): vielbein postulate d\_mu e\textasciicircum{}A\_nu - chr\textasciicircum{}l\_mu nu e\textasciicircum{}A\_l + omega\_mu\textasciicircum{}A\_B e\textasciicircum{}B\_nu = 0} & True \\
65 & \texttt{inhomogeneous metric N, A, B (x0, x4): Omega\_mu real, Gamma17[0] Omega\_mu\textasciicircum{}T Gamma17[0] = -Omega\_mu} & True \\
66 & \texttt{inhomogeneous metric N, A, B (x0, x4): total covariant derivative of the frame matrices: d\_mu Gamma\_nu - chr\textasciicircum{}l\_mu nu Gamma\_l + [Omega\_mu, Gamma\_nu] = 0} & True \\
67 & \texttt{inhomogeneous metric N, A, B (x0, x4): constant frame matrices: omega\_mu\textasciicircum{}A\_B Gamma\textasciicircum{}B + [Omega\_mu, Gamma\textasciicircum{}A] = 0} & True \\
68 & \texttt{inhomogeneous metric N, A, B (x0, x4): spin curvature F\_mu nu = (1/4) R\_AB mu nu Gamma\textasciicircum{}A Gamma\textasciicircum{}B} & True \\
69 & \texttt{inhomogeneous metric N, A, B (x0, x4): Bianchi identity d\_l F\_mn + [Omega\_l, F\_mn] + cyclic = 0} & True \\
70 & \texttt{inhomogeneous metric N, A, B (x0, x4): (1/2) Gamma\textasciicircum{}mu Gamma\textasciicircum{}nu F\_mu nu = -(R/4) 1\_16} & True \\
71 & \texttt{Cartan's first structure equation de\textasciicircum{}A + omega\textasciicircum{}A\_B \textasciicircum{} e\textasciicircum{}B = 0 (torsion-free) for the cosmological frame} & True \\
72 & \texttt{static radial metric f(r) dx0\textasciicircum{}2 - h(r) dr\textasciicircum{}2 - r\textasciicircum{}2 (dx2\textasciicircum{}2 + ... + dx7\textasciicircum{}2): eta e e = g} & True \\
73 & \texttt{static radial metric f(r) dx0\textasciicircum{}2 - h(r) dr\textasciicircum{}2 - r\textasciicircum{}2 (dx2\textasciicircum{}2 + ... + dx7\textasciicircum{}2): \{Gamma\_mu, Gamma\_nu\} = 2 g\_mu nu} & True \\
74 & \texttt{static radial metric f(r) dx0\textasciicircum{}2 - h(r) dr\textasciicircum{}2 - r\textasciicircum{}2 (dx2\textasciicircum{}2 + ... + dx7\textasciicircum{}2): Gamma\_mu real} & True \\
75 & \texttt{static radial metric f(r) dx0\textasciicircum{}2 - h(r) dr\textasciicircum{}2 - r\textasciicircum{}2 (dx2\textasciicircum{}2 + ... + dx7\textasciicircum{}2): omega\_mu AB = -omega\_mu BA} & True \\
76 & \texttt{static radial metric f(r) dx0\textasciicircum{}2 - h(r) dr\textasciicircum{}2 - r\textasciicircum{}2 (dx2\textasciicircum{}2 + ... + dx7\textasciicircum{}2): vielbein postulate d\_mu e\textasciicircum{}A\_nu - chr\textasciicircum{}l\_mu nu e\textasciicircum{}A\_l + omega\_mu\textasciicircum{}A\_B e\textasciicircum{}B\_nu = 0} & True \\
77 & \texttt{static radial metric f(r) dx0\textasciicircum{}2 - h(r) dr\textasciicircum{}2 - r\textasciicircum{}2 (dx2\textasciicircum{}2 + ... + dx7\textasciicircum{}2): Omega\_mu real, Gamma17[0] Omega\_mu\textasciicircum{}T Gamma17[0] = -Omega\_mu} & True \\
78 & \texttt{static radial metric f(r) dx0\textasciicircum{}2 - h(r) dr\textasciicircum{}2 - r\textasciicircum{}2 (dx2\textasciicircum{}2 + ... + dx7\textasciicircum{}2): total covariant derivative of the frame matrices: d\_mu Gamma\_nu - chr\textasciicircum{}l\_mu nu Gamma\_l + [Omega\_mu, Gamma\_nu] = 0} & True \\
79 & \texttt{static radial metric f(r) dx0\textasciicircum{}2 - h(r) dr\textasciicircum{}2 - r\textasciicircum{}2 (dx2\textasciicircum{}2 + ... + dx7\textasciicircum{}2): constant frame matrices: omega\_mu\textasciicircum{}A\_B Gamma\textasciicircum{}B + [Omega\_mu, Gamma\textasciicircum{}A] = 0} & True \\
80 & \texttt{static radial metric f(r) dx0\textasciicircum{}2 - h(r) dr\textasciicircum{}2 - r\textasciicircum{}2 (dx2\textasciicircum{}2 + ... + dx7\textasciicircum{}2): spin curvature F\_mu nu = (1/4) R\_AB mu nu Gamma\textasciicircum{}A Gamma\textasciicircum{}B} & True \\
81 & \texttt{static radial metric f(r) dx0\textasciicircum{}2 - h(r) dr\textasciicircum{}2 - r\textasciicircum{}2 (dx2\textasciicircum{}2 + ... + dx7\textasciicircum{}2): Bianchi identity d\_l F\_mn + [Omega\_l, F\_mn] + cyclic = 0} & True \\
82 & \texttt{static radial metric f(r) dx0\textasciicircum{}2 - h(r) dr\textasciicircum{}2 - r\textasciicircum{}2 (dx2\textasciicircum{}2 + ... + dx7\textasciicircum{}2): (1/2) Gamma\textasciicircum{}mu Gamma\textasciicircum{}nu F\_mu nu = -(R/4) 1\_16} & True \\
83 & \texttt{Cartan's first structure equation for the static radial metric} & True \\
84 & \texttt{structure coefficients: de\textasciicircum{}A = (1/2) C\textasciicircum{}A\_BC e\textasciicircum{}B \textasciicircum{} e\textasciicircum{}C and omega\_C,AB (form index C) = -(1/2)(C\_ACB + C\_BAC - C\_CBA) (static radial metric)} & True \\
85 & \texttt{static radial metric: Omega\_0 = (f'/(4 sqrt(f h))) Gamma17[1] Gamma17[0], Omega\_1 = 0, Omega\_k = (1/(2 sqrt h)) Gamma17[1] Gamma17[k] (k = 2..7)} & True \\
86 & \texttt{static radial metric: Gamma\textasciicircum{}mu Omega\_mu = (1/sqrt h)(f'/(4f) + 3/r) Gamma17up[1], so (i Gamma\textasciicircum{}mu D\_mu - m) psi = [i Gamma17up[0] d0/sqrt f + i Gamma17up[1] (d\_r + f'/(4f) + 3/r)/sqrt h + sum\_k i Gamma17up[k] d\_k/r - m] psi} & True \\
87 & \texttt{static radial metric: psi = f\textasciicircum{}(-1/4) r\textasciicircum{}(-3) chi removes the spin connection: Gamma\textasciicircum{}mu D\_mu psi = f\textasciicircum{}(-1/4) r\textasciicircum{}(-3) Gamma\textasciicircum{}mu d\_mu chi} & True \\
88 & \texttt{U(1) gauge covariance: Gamma\textasciicircum{}mu (D\_mu - i q A'\_mu)(e\textasciicircum{}\{i q alpha\} psi) = e\textasciicircum{}\{i q alpha\} Gamma\textasciicircum{}mu (D\_mu - i q A\_mu) psi with A' = A + d alpha} & True \\
89 & \texttt{current identity with the gauge field: d\_mu(sqrt|g| psibar Gamma\textasciicircum{}mu psi) = sqrt|g| [(D\_mu psibar) Gamma\textasciicircum{}mu psi + psibar Gamma\textasciicircum{}mu D\_mu psi], D including -i q A\_mu} & True \\
90 & \texttt{current identity d\_mu(sqrt|g| psibar Gamma\textasciicircum{}mu psi) = sqrt|g| [(D\_mu psibar) Gamma\textasciicircum{}mu psi + psibar Gamma\textasciicircum{}mu D\_mu psi]} & True \\
91 & \texttt{Lichnerowicz: (Gamma\textasciicircum{}mu D\_mu)\textasciicircum{}2 psi = (g\textasciicircum{}mu nu nabla\_mu D\_nu - R/4) psi} & True \\
92 & \texttt{rescaling psi = (a\textasciicircum{}3 b\textasciicircum{}4)\textasciicircum{}(-1/2) chi removes the spin connection: Gamma\textasciicircum{}mu D\_mu psi = (a\textasciicircum{}3 b\textasciicircum{}4)\textasciicircum{}(-1/2) Gamma\textasciicircum{}mu d\_mu chi} & True \\
93 & \texttt{mode equation: (i Gamma\textasciicircum{}mu D\_mu - m) psi = 0  <=>  i dc/dx0 = h(x0) c, h = alpha.(p/a, q/b) + beta m} & True \\
94 & \texttt{conformal weight 7/2: Gamma\textasciicircum{}mu D\_mu (C\textasciicircum{}(-7/2) f) = C\textasciicircum{}(-9/2) Gamma\textasciicircum{}A d\_A f for g = C(x0)\textasciicircum{}2 eta} & True \\
95 & \texttt{2x2 reduction: A = alpha.n\_a, B = alpha.n\_b, beta: squares 1, mutually anticommuting; Z = A B beta: Z\textasciicircum{}2 = -1, Z\textasciicircum{}T = -Z, commutes with A, B, beta} & True \\
96 & \texttt{Kasner: |alpha|\textasciicircum{}2 + |beta|\textasciicircum{}2 = 1 (to 1e-8)} & True \\
97 & \texttt{Kasner scaling: (|p|, |q|) = (1.430..., 2.5) has K = 1 and the same |beta|\textasciicircum{}2 as (1, 1) (to 1e-6)} & True \\
98 & \texttt{16x16 evolution = 4 copies of the 2x2 evolution with +m sigma3 and 4 copies with -m sigma3: average agrees (to 1e-8), unitarity; the two sectors differ} & True \\
99 & \texttt{isotropic expansion, massless: no pair creation (|beta|\textasciicircum{}2 < 1e-20)} & True \\
100 & \texttt{de Sitter: |beta|\textasciicircum{}2 = 1/(exp(2 pi m/H) + 1) (Fermi-Dirac at the Gibbons-Hawking temperature H/2pi), to 1e-7} & True \\
101 & \texttt{integrator benchmark: Landau-Zener |beta|\textasciicircum{}2 = exp(-pi kappa\textasciicircum{}2) (to 1e-6)} & True \\*
\end{longtable}}
\textbf{Numbers computed by both programs:}

{\small
\begin{longtable}{>{\raggedright\arraybackslash}p{\dimexpr 0.5000\linewidth-2\tabcolsep\relax}>{\raggedright\arraybackslash}p{\dimexpr 0.2500\linewidth-2\tabcolsep\relax}>{\raggedright\arraybackslash}p{\dimexpr 0.2500\linewidth-2\tabcolsep\relax}}
\toprule
\textbf{quantity} & \textbf{7+1.nb} & \textbf{Python} \\
\midrule\endfirsthead
\toprule
\textbf{quantity} & \textbf{7+1.nb} & \textbf{Python} \\
\midrule\endhead
\endfoot
\bottomrule\endlastfoot
Kasner $|\beta|^2$ at $(|\vec p|,|\vec q|)=(1,1)$ & 0.010749346 & 0.010749346 \\
Kasner $|\beta|^2$ at $(|\vec p|,|\vec q|)=(2,1)$ & 0.00013294108 & 0.00013294108 \\
Kasner $|\beta|^2$ at $(|\vec p|,|\vec q|)=(1,3)$ & 0.067442022 & 0.067442022 \\
de Sitter $|\beta|^2$, $M/H=0.25$ & 0.172102919 & — \\
de Sitter $|\beta|^2$, $M/H=0.5$ & 0.0414238845 & — \\
de Sitter $|\beta|^2$, $M/H=1$ & 0.00186397746 & — \\*
\end{longtable}}
\textbf{SHA-256 of the files} (at the time this document was generated):

\begin{Verbatim}[fontsize=\small,frame=single,framesep=2mm]
7+1.nb                             1314b1e09e6b6eefbd9ee5974e19d186a5e996907b2eb68ddb6f23aae85aa
    ...aa5
tools/make_7p1_nb.py               354a0210e60510499753c2cf91274850cedef502d6f705efc410c5a8bbae7
    ...b2d
tools/nb_part_gravity.py           c4caf2ad5c932e8f603341edade5f4b98fa784feb05906c051eab7c884cd6
    ...8f8
tools/nb_part_ks_gravity.py        a656c0ad040b4576b75b17e029545109e531c8e656ad99bfdb5cba946e967
    ...772
tools/nb_part_sm.py                43e92beed57453a4b57931de6d603ac6c218efba6e4005133b69b8054a975
    ...c5b
tools/nb_part_coupled.py           49c61aeb6affe14ebbcccb32d2069504b9fc035313ea2fabbd0a3392a9613
    ...3f5
tools/lovelock71.py                d93ecbd0f82adb972218d59dfe32999a4af944945ca9758607f3a633e1af9
    ...5a1
tools/ref_gravity.py               f07cf8adaff585af89c87d19532f1aaa7db606ad88ae90483b2d1f41341ed
    ...f36
tools/ref_ks_gravity.py            b6488ef24b1c16aacd8ff6cf6a5a1ba3f316e082449ad97480022d1a26383
    ...27f
tools/ref_sm.py                    f3f80806893f42bbf166e1f1de896ffd611dbfe847685da8e3ee367e9cde6
    ...aba
tools/ref_coupled.py               695dbd9375753653472d9c37502f0067c4a64fe01a64ba8c05e8b2c5da85a
    ...a58
tools/ref71.py                     1a7a06eb0add3f459611776dc9c474a5bb959d947a9db949db4612b1ab943
    ...4e0
tools/ref44.py                     ab100c3ebcc055fbe59c132a02f80ecb17630e934e99d59f88e919ec8064f
    ...af5
\end{Verbatim}
\textbf{Output of tools/ref\_gravity.py (independent checks of sections 11-14; exit code 0)} (complete):

\begin{Verbatim}[fontsize=\small,frame=single,framesep=2mm]
eta17 is simply the squares of the octonion basis: eta_AB = Re(e_A e_B), e_0^2 = +1, e_i^2 = -1
    ...(i = 1..7) : True
the quadratic form of the octonions has the signs of the squares of the octonion basis: Re(x x)
    ...= x0^2 - x1^2 - ... - x7^2 = x.eta17.x : True
cosmological metric: sectional curvatures R^{AB}_{AB} = -(Ha'+Ha^2), -(Hb'+Hb^2), -Ha^2, -Hb^2,
    ...-Ha Hb : True
cosmological metric: Riemann R^{cd}_{ab} non-zero only for {c,d} = {a,b} : True
G^0_0 = 3Ha^2 + 12 Ha Hb + 6 Hb^2 : True
G^1_1 = 2Ha' + 3Ha^2 + 4Hb' + 10Hb^2 + 8HaHb : True
G^4_4 = 3Ha' + 6Ha^2 + 3Hb' + 6Hb^2 + 9HaHb : True
scalar curvature R = -2(3Ha' + 6Ha^2 + 4Hb' + 10Hb^2 + 12 HaHb) : True
E^(1) = Einstein tensor G (cosmological metric) : True
L_1 = R : True
E^(0) = -delta/2 : True
trace identity E^(k)mu_mu = -(8-2k)/2 L_k, k = 0..4 : True
E^(4) = 0 identically in 8 dimensions, but L_4 (Euler density) is not 0 : True
E^(k) diagonal: E^(k)mu_nu = 0 for mu != nu, and E^(k)1_1 = E^(k)2_2 = E^(k)3_3, E^(k)4_4 = ...
    ...= E^(k)7_7 : True
E^(2)0_0 = -12 Hb (4Ha^3 + 18Ha^2Hb + 12HaHb^2 + Hb^3) : True
E^(3)0_0 = 360 Ha^2 Hb^3 (4Ha + 3Hb) : True
E^(2) = Lanczos tensor 2(R R_mn - 2R_ma R^a_n - 2R^ab R_manb + R_m^abc R_nabc) - g_mn L_2/2,
    ...and L_2 = R^2 - 4 Ric.Ric + Riem.Riem (cosmological metric) : True
closed forms : {(1, 0): 3*H_a**2 + 12*H_a*H_b + 6*H_b**2, (1, 1): 3*H_a**2 + 8*H_a*H_b +
    ...10*H_b**2 + 2*dH_a + 4*dH_b, (1, 4): 6*H_a**2 + 9*H_a*H_b + 6*H_b**2 + 3*dH_a + 3*dH_b,
    ...(2, 0): -48*H_a**3*H_b - 216*H_a**2*H_b**2 - 144*H_a*H_b**3 - 12*H_b**4, (2, 1):
    ...-32*H_a**3*H_b - 136*H_a**2*H_b**2 - 16*H_a**2*dH_b - 192*H_a*H_b**3 - 32*H_a*H_b*dH_a -
    ...96*H_a*H_b*dH_b - 60*H_b**4 - 48*H_b**2*dH_a - 48*H_b**2*dH_b, (2, 4): -12*H_a**4 -
    ...108*H_a**3*H_b - 180*H_a**2*H_b**2 - 12*H_a**2*dH_a - 36*H_a**2*dH_b - 108*H_a*H_b**3 -
    ...72*H_a*H_b*dH_a - 72*H_a*H_b*dH_b - 12*H_b**4 - 36*H_b**2*dH_a - 12*H_b**2*dH_b, (3, 0):
    ...1440*H_a**3*H_b**3 + 1080*H_a**2*H_b**4, (3, 1): 576*H_a**3*H_b**3 + 1368*H_a**2*H_b**4
    ...+ 864*H_a**2*H_b**2*dH_b + 576*H_a*H_b**5 + 576*H_a*H_b**3*dH_a + 576*H_a*H_b**3*dH_b +
    ...144*H_b**4*dH_a, (3, 4): 648*H_a**4*H_b**2 + 1224*H_a**3*H_b**3 + 432*H_a**3*H_b*dH_b +
    ...648*H_a**2*H_b**4 + 648*H_a**2*H_b**2*dH_a + 648*H_a**2*H_b**2*dH_b +
    ...432*H_a*H_b**3*dH_a}
divergence-free E^(k), k = 1, 2, 3 (cosmological metric, arbitrary a, b) : True
variation of sqrt|g| L_k (lapse N, a, b): EL_N = -2 a^3 b^4 E^(k)0_0, EL_a = -6 a^2 b^4
    ...E^(k)1_1, EL_b = -8 a^3 b^3 E^(k)4_4, k = 1..4 : True
inhomogeneous metric: E^(3)0_4 at the point (exact) :
    ...35374305025486636433298941367279331645321641984000/24577736888913221711303629009057562607
    ...63178144678501
inhomogeneous metric (N, A, B of x0 and x4): E^(1) = G at the point : True
inhomogeneous metric: nabla_mu E^(k)mu_nu = 0 exactly at x0 = 1/2, x4 = 1/3, k = 1, 2, 3 : True
inhomogeneous metric: E^(k)_mu nu symmetric, k = 1, 2, 3 : True
pp-wave: E^(1)_mu nu = +(1/2) (d2^2 + d3^2) F l_mu l_nu, l = dx0 - dx1 : True
pp-wave: E^(2) = E^(3) = 0 for every profile F : True
Kasner E^(k) components x0^(2k) (00, 11, 44) : {1: [3*(p**2 + 4*p*q + 2*q**2), 3*p**2 + 8*p*q -
    ...2*p + 10*q**2 - 4*q, 3*(2*p**2 + 3*p*q - p + 2*q**2 - q)], 2: [-12*q*(4*p**3 + 18*p**2*q
    ...+ 12*p*q**2 + q**3), -4*q*(8*p**3 + 34*p**2*q - 12*p**2 + 48*p*q**2 - 36*p*q + 15*q**3 -
    ...12*q**2), -12*(p**4 + 9*p**3*q - p**3 + 15*p**2*q**2 - 9*p**2*q + 9*p*q**3 - 9*p*q**2 +
    ...q**4 - q**3)], 3: [360*p**2*q**3*(4*p + 3*q), 72*p*q**3*(8*p**2 + 19*p*q - 20*p + 8*q**2
    ...- 10*q), 72*p**2*q**2*(9*p**2 + 17*p*q - 15*p + 9*q**2 - 15*q)]}
Einstein Kasner: p = (1+-2 sqrt2)/7, q = (2-+3 sqrt2)/14 solve E^(1) = 0; 3p+4q = 1, 3p^2+4q^2
    ...= 1 : True
Gauss-Bonnet Kasner cubic (64 x e4 at q = (3-3p)/4) : -203*p**3 + 81*p**2 + 351*p + 27
Gauss-Bonnet Kasner exponents p (three real roots) : (-1.08350178761210574242397629473,
    ...-0.0786310584189277662278610521099, 1.56114762435615666136119695275)
Gauss-Bonnet Kasner: the three (p, (3-3p)/4) solve E^(2) = 0, sum of the 7 exponents = 3 : True
cubic Lovelock Kasner: (p, q) = (-15/7, 20/7) solves E^(3) = 0, 3p + 4q = 5; degenerate
    ...families q = 0 and p = 0 : True
Einstein Kasner is not a cubic-Lovelock vacuum (E^(3) != 0) : True
8D de Sitter (a = b = e^{H x0}): E^(1) = 21 H^2, E^(2) = -420 H^4, E^(3) = 2520 H^6 (times
    ...delta) : True
4D de Sitter on a static T^4 (a = e^{H x0}, b = b0): E^(1) = diag(3,3,6) H^2, E^(2) =
    ...diag(0,0,-12) H^4, E^(3) = 0 : True
static T^4 (Hb = 0): E^(1) = (3Ha^2, 2Ha'+3Ha^2, 3Ha'+6Ha^2), E^(2) = (0, 0, -12 Ha^2
    ...(Ha'+Ha^2)), E^(3) = 0 : True
exponential expansion a = e^{H x0}, b = e^{q H x0}: E^(k) (00, 11, 44) = closed forms with Ha =
    ...H, Hb = q H, Ha' = Hb' = 0 (k = 1, 2, 3) : True
exponential expansion: 3 Ha (eq00 - eq11) + 4 Hb (eq00 - eq44) = 0 identically (the three
    ...vacuum equations have rank 2) : True
exponential vacuum solutions: one solution (w2, w3) for given w1, Lambda, q, satisfying all
    ...three equations : True
eternal inflation a = e^{H x0} with deflation b = e^{-H x0}: w2 = Lambda/4 - w1/2, w3 = w1/24 -
    ...Lambda/36 (Lambda = 0: w2/w1 = -1/2, w3/w1 = 1/24) : True
NUM exponential vacuum solution (w2, w3) for general q : ((-2*Lambda*q - 3*Lambda + 12*q**3*w1
    ...+ 27*q**2*w1 + 22*q*w1 + 9*w1)/(4*q*(6*q**4 + 36*q**3 + 66*q**2 + 51*q + 16)),
    ...(-3*Lambda*q**2 - 6*Lambda*q - Lambda + 15*q**4*w1 + 36*q**3*w1 + 33*q**2*w1 + 18*q*w1 +
    ...3*w1)/(72*q**3*(6*q**4 + 36*q**3 + 66*q**2 + 51*q + 16)))
NUM eq00, eq11, eq44 for q = -1 : (-Lambda - 3*w1 - 36*w2 - 360*w3, -Lambda + 5*w1 + 28*w2 +
    ...216*w3, -Lambda + 3*w1 + 12*w2 + 72*w3)
cosmological metric: eta e e = g : True
cosmological metric: {Gamma_mu, Gamma_nu} = 2 g_mu nu : True
cosmological metric: Gamma_mu real : True
cosmological metric: omega_mu AB = -omega_mu BA : True
cosmological metric: vielbein postulate d_mu e^A_nu - chr^l_mu nu e^A_l + omega_mu^A_B e^B_nu =
    ...0 : True
cosmological metric: Omega_mu real, Gamma17[0] Omega_mu^T Gamma17[0] = -Omega_mu : True
cosmological metric: total covariant derivative of the frame matrices: d_mu Gamma_nu - chr^l_mu
    ...nu Gamma_l + [Omega_mu, Gamma_nu] = 0 : True
cosmological metric: constant frame matrices: omega_mu^A_B Gamma^B + [Omega_mu, Gamma^A] = 0 :
    ...True
cosmological metric: spin curvature F_mu nu = (1/4) R_AB mu nu Gamma^A Gamma^B : True
cosmological metric: Bianchi identity d_l F_mn + [Omega_l, F_mn] + cyclic = 0 : True
cosmological metric: (1/2) Gamma^mu Gamma^nu F_mu nu = -(R/4) 1_16 : True
cosmological metric: Gamma^mu Omega_mu = (3a'/(2a) + 2b'/b) Gamma17[0] : True
cosmological metric: Omega_i = (a'/2) Gamma17[i] Gamma17[0] (i = 1..3), Omega_k = (b'/2)
    ...Gamma17[k] Gamma17[0] (k = 4..7), Omega_0 = 0 : True
pp-wave: eta e e = g : True
pp-wave: {Gamma_mu, Gamma_nu} = 2 g_mu nu : True
pp-wave: Gamma_mu real : True
pp-wave: omega_mu AB = -omega_mu BA : True
pp-wave: vielbein postulate d_mu e^A_nu - chr^l_mu nu e^A_l + omega_mu^A_B e^B_nu = 0 : True
pp-wave: Omega_mu real, Gamma17[0] Omega_mu^T Gamma17[0] = -Omega_mu : True
pp-wave: total covariant derivative of the frame matrices: d_mu Gamma_nu - chr^l_mu nu Gamma_l
    ...+ [Omega_mu, Gamma_nu] = 0 : True
pp-wave: constant frame matrices: omega_mu^A_B Gamma^B + [Omega_mu, Gamma^A] = 0 : True
pp-wave: spin curvature F_mu nu = (1/4) R_AB mu nu Gamma^A Gamma^B : True
pp-wave: Bianchi identity d_l F_mn + [Omega_l, F_mn] + cyclic = 0 : True
pp-wave: (1/2) Gamma^mu Gamma^nu F_mu nu = -(R/4) 1_16 : True
inhomogeneous metric N, A, B (x0, x4): eta e e = g : True
inhomogeneous metric N, A, B (x0, x4): {Gamma_mu, Gamma_nu} = 2 g_mu nu : True
inhomogeneous metric N, A, B (x0, x4): Gamma_mu real : True
inhomogeneous metric N, A, B (x0, x4): omega_mu AB = -omega_mu BA : True
inhomogeneous metric N, A, B (x0, x4): vielbein postulate d_mu e^A_nu - chr^l_mu nu e^A_l +
    ...omega_mu^A_B e^B_nu = 0 : True
inhomogeneous metric N, A, B (x0, x4): Omega_mu real, Gamma17[0] Omega_mu^T Gamma17[0] =
    ...-Omega_mu : True
inhomogeneous metric N, A, B (x0, x4): total covariant derivative of the frame matrices: d_mu
    ...Gamma_nu - chr^l_mu nu Gamma_l + [Omega_mu, Gamma_nu] = 0 : True
inhomogeneous metric N, A, B (x0, x4): constant frame matrices: omega_mu^A_B Gamma^B +
    ...[Omega_mu, Gamma^A] = 0 : True
inhomogeneous metric N, A, B (x0, x4): spin curvature F_mu nu = (1/4) R_AB mu nu Gamma^A
    ...Gamma^B : True
inhomogeneous metric N, A, B (x0, x4): Bianchi identity d_l F_mn + [Omega_l, F_mn] + cyclic = 0
    ...: True
inhomogeneous metric N, A, B (x0, x4): (1/2) Gamma^mu Gamma^nu F_mu nu = -(R/4) 1_16 : True
Cartan's first structure equation de^A + omega^A_B ^ e^B = 0 (torsion-free) for the
    ...cosmological frame : True
static radial metric f(r) dx0^2 - h(r) dr^2 - r^2 (dx2^2 + ... + dx7^2): eta e e = g : True
static radial metric f(r) dx0^2 - h(r) dr^2 - r^2 (dx2^2 + ... + dx7^2): {Gamma_mu, Gamma_nu} =
    ...2 g_mu nu : True
static radial metric f(r) dx0^2 - h(r) dr^2 - r^2 (dx2^2 + ... + dx7^2): Gamma_mu real : True
static radial metric f(r) dx0^2 - h(r) dr^2 - r^2 (dx2^2 + ... + dx7^2): omega_mu AB =
    ...-omega_mu BA : True
static radial metric f(r) dx0^2 - h(r) dr^2 - r^2 (dx2^2 + ... + dx7^2): vielbein postulate
    ...d_mu e^A_nu - chr^l_mu nu e^A_l + omega_mu^A_B e^B_nu = 0 : True
static radial metric f(r) dx0^2 - h(r) dr^2 - r^2 (dx2^2 + ... + dx7^2): Omega_mu real,
    ...Gamma17[0] Omega_mu^T Gamma17[0] = -Omega_mu : True
static radial metric f(r) dx0^2 - h(r) dr^2 - r^2 (dx2^2 + ... + dx7^2): total covariant
    ...derivative of the frame matrices: d_mu Gamma_nu - chr^l_mu nu Gamma_l + [Omega_mu,
    ...Gamma_nu] = 0 : True
static radial metric f(r) dx0^2 - h(r) dr^2 - r^2 (dx2^2 + ... + dx7^2): constant frame
    ...matrices: omega_mu^A_B Gamma^B + [Omega_mu, Gamma^A] = 0 : True
static radial metric f(r) dx0^2 - h(r) dr^2 - r^2 (dx2^2 + ... + dx7^2): spin curvature F_mu nu
    ...= (1/4) R_AB mu nu Gamma^A Gamma^B : True
static radial metric f(r) dx0^2 - h(r) dr^2 - r^2 (dx2^2 + ... + dx7^2): Bianchi identity d_l
    ...F_mn + [Omega_l, F_mn] + cyclic = 0 : True
static radial metric f(r) dx0^2 - h(r) dr^2 - r^2 (dx2^2 + ... + dx7^2): (1/2) Gamma^mu
    ...Gamma^nu F_mu nu = -(R/4) 1_16 : True
Cartan's first structure equation for the static radial metric : True
structure coefficients: de^A = (1/2) C^A_BC e^B ^ e^C and omega_C,AB (form index C) =
    ...-(1/2)(C_ACB + C_BAC - C_CBA) (static radial metric) : True
static radial metric: Omega_0 = (f'/(4 sqrt(f h))) Gamma17[1] Gamma17[0], Omega_1 = 0, Omega_k
    ...= (1/(2 sqrt h)) Gamma17[1] Gamma17[k] (k = 2..7) : True
static radial metric: Gamma^mu Omega_mu = (1/sqrt h)(f'/(4f) + 3/r) Gamma17up[1], so (i
    ...Gamma^mu D_mu - m) psi = [i Gamma17up[0] d0/sqrt f + i Gamma17up[1] (d_r + f'/(4f) +
    ...3/r)/sqrt h + sum_k i Gamma17up[k] d_k/r - m] psi : True
static radial metric: psi = f^(-1/4) r^(-3) chi removes the spin connection: Gamma^mu D_mu psi
    ...= f^(-1/4) r^(-3) Gamma^mu d_mu chi : True
U(1) gauge covariance: Gamma^mu (D_mu - i q A'_mu)(e^{i q alpha} psi) = e^{i q alpha} Gamma^mu
    ...(D_mu - i q A_mu) psi with A' = A + d alpha : True
current identity with the gauge field: d_mu(sqrt|g| psibar Gamma^mu psi) = sqrt|g| [(D_mu
    ...psibar) Gamma^mu psi + psibar Gamma^mu D_mu psi], D including -i q A_mu : True
current identity d_mu(sqrt|g| psibar Gamma^mu psi) = sqrt|g| [(D_mu psibar) Gamma^mu psi +
    ...psibar Gamma^mu D_mu psi] : True
Lichnerowicz: (Gamma^mu D_mu)^2 psi = (g^mu nu nabla_mu D_nu - R/4) psi : True
rescaling psi = (a^3 b^4)^(-1/2) chi removes the spin connection: Gamma^mu D_mu psi = (a^3
    ...b^4)^(-1/2) Gamma^mu d_mu chi : True
mode equation: (i Gamma^mu D_mu - m) psi = 0  <=>  i dc/dx0 = h(x0) c, h = alpha.(p/a, q/b) +
    ...beta m : True
conformal weight 7/2: Gamma^mu D_mu (C^(-7/2) f) = C^(-9/2) Gamma^A d_A f for g = C(x0)^2 eta :
    ...True
2x2 reduction: A = alpha.n_a, B = alpha.n_b, beta: squares 1, mutually anticommuting; Z = A B
    ...beta: Z^2 = -1, Z^T = -Z, commutes with A, B, beta : True
Kasner (Einstein branch), massless: |beta|^2 for (|p|, |q|) = (1,1), (2,1), (1,3) :
    ...(np.float64(0.01074934616777871), np.float64(0.00013294108220674155),
    ...np.float64(0.06744202228393215))
Kasner scaling variable K = |p| lambda^(1-p), lambda = (|q|/|p|)^(1/(q-p)) : (1.0,
    ...3.11828865465698, 0.4946322370571032)
Kasner: |alpha|^2 + |beta|^2 = 1 (to 1e-8) : True
Kasner scaling: (|p|, |q|) = (1.430..., 2.5) has K = 1 and the same |beta|^2 as (1, 1) (to
    ...1e-6) : True
Kasner, m = 0.3, |p| = 1, |q| = 1.5, x0: 0.05 -> 40: |beta|^2 of the 16x16 evolution (average),
    ...of the sector Z = +i (+m sigma3) and of the sector Z = -i (-m sigma3) :
    ...(np.float64(0.027052647156543617), np.float64(0.02528659275098409),
    ...np.float64(0.028818701562187667))
16x16 evolution = 4 copies of the 2x2 evolution with +m sigma3 and 4 copies with -m sigma3:
    ...average agrees (to 1e-8), unitarity; the two sectors differ : True
isotropic expansion, massless: no pair creation (|beta|^2 < 1e-20) : True
far future (de Sitter, H = 1): |beta|^2 for m/H = 0.25, 0.5, 1 :
    ...(np.float64(0.1721029191835808), np.float64(0.04142388447252214),
    ...np.float64(0.001863977463333981))
de Sitter: |beta|^2 = 1/(exp(2 pi m/H) + 1) (Fermi-Dirac at the Gibbons-Hawking temperature
    ...H/2pi), to 1e-7 : True
integrator benchmark: Landau-Zener |beta|^2 = exp(-pi kappa^2) (to 1e-6) : True
boolean checks: 101   failed: 0
VERDICT: ALL CHECKS PASSED
\end{Verbatim}
\end{document}
